\documentclass[11pt,a4paper]{article}
\usepackage[OT1]{fontenc}
\usepackage{amsmath,amssymb,amsthm,mathtools}
\usepackage[margin=27mm]{geometry}
\usepackage{microtype,booktabs,longtable,array,enumitem}
\usepackage[hidelinks]{hyperref}
\setlength{\emergencystretch}{2em}
\setlength{\parskip}{3pt plus 1pt}
\setcounter{tocdepth}{1}
\numberwithin{equation}{section}
\newtheorem{theorem}{Theorem}[section]
\newtheorem{lemma}[theorem]{Lemma}
\newtheorem{proposition}[theorem]{Proposition}
\newtheorem{corollary}[theorem]{Corollary}
\theoremstyle{definition}
\newtheorem{definition}[theorem]{Definition}
\newtheorem{example}[theorem]{Example}
\theoremstyle{remark}
\newtheorem{remark}[theorem]{Remark}
\newcommand{\N}{\mathbb N}
\newcommand{\NN}{\mathbb N_{>0}}
\newcommand{\norm}{N}
\newcommand{\can}[1]{\mathopen{[\![}#1\mathclose{]\!]}}
\newcommand{\mle}{\preccurlyeq}
\newcommand{\mslt}{\prec}
\newcommand{\lv}{\operatorname{lv}}
\newcommand{\sz}{\operatorname{sz}}
\newcommand{\FV}{\operatorname{FV}}
\newcommand{\Ty}{\operatorname{Ty}}
\title{Decidable Ground Comparison for\\Norm-Controlled Ordinal Assignments}
\author{}
\date{17 November 2026\\Research manuscript}
\begin{document}
\maketitle
\begin{abstract}
Wilken and Weiermann's ordinal assignment for G\"odel's T uses a rapidly
growing primitive recursive control and comparison under all admissible
substitutions. We give explicit permitted controls for which universal
comparison of canonical assignments is decidable whenever all free term
variables have ground type, without restrictions on internal types,
recursors, or finite nesting. One control preserves any prescribed finite
prefix of the published example and changes its later values by less than
twice the input factorial. The main arithmetic observation is a modular
potential that decreases by one at every norm-controlled descent event.
A final passage through \(\omega^2\) then determines complete descent
counts modulo every integer at every fixed ordinal head. We prove the
entry comparison and finite elimination arguments needed to turn these
tables into a decision procedure. For the same controls, weak and strict
comparison become undecidable with free variables of type \(o\to o\).
The proofs use a literal canonical assignment and corrected scalar and
norm estimates. The all-ground problem for the unchanged published
control is not decided here.
\end{abstract}
\tableofcontents

\section{The result and its scope}\label{sec:introduction}

A termination argument need not give an algorithm for comparing the
measures it assigns. In the analysis of G\"odel's system T by Wilken and
Weiermann \cite{ww}, ordinal expressions are compared under every
admissible substitution for their variables. Definition~2.8 of that
paper asks whether these relations can be made effective. Its collapse
construction permits a choice of sufficiently fast primitive recursive
control, subject to an eventual bound at a fixed finite hierarchy level.
The displayed example is \(G(n)=F_5(n+100)\), with
\(F_0(n)=2^n\) and \(F_{j+1}(n)=F_j^{n+1}(n)\).

We fix the \emph{canonical} assignment, which uses a history of length
zero at each abstraction. Its literal syntax matters: rebracketing an
existing coordinate expression can change a later syntactic partial
operator. A computable admissible parsing policy therefore forms part of
the definition. It preserves the input subtrees at each macro site and
commutes with substitution. Section~\ref{sec:source} specifies that
definition and the repaired inequalities it needs.

For vectors \(a,b\), weak comparison means universal coordinatewise
inequality under all admissible valuations. Wilken--Weiermann strict
comparison means universal strict inequality in coordinate zero and
universal weak inequality in every positive coordinate. Equality is
universal coordinatewise equality. These are comparisons of assigned
vectors, not beta equivalence of the underlying terms.

\begin{theorem}[Comparison classification]\label{thm:main}
There is an explicit permitted primitive recursive control \(f_*\) for
which the following hold, uniformly in a supplied computable admissible
literal parsing policy.
\begin{enumerate}[label=(\roman*)]
\item Universal canonical weak comparison, strict comparison and equality
are decidable for original T terms all of whose free variables have ground
type. Bound variables, result types, internal types, and recursors are
unrestricted.
\item Universal canonical weak and strict comparison are undecidable for
ground-output terms with free-variable levels at most one and internal
type levels at most two, by a finite-query reduction.
\item The negative result has a single-query reduction when internal
type levels at most three are allowed.
\end{enumerate}
For every supplied \(b\), the same assertions hold for an explicit
control \(f_b\) equal to \(G\) on \([0,b]\) and satisfying
\(0\le f_b(n)-G(n)<2\operatorname{lcm}(1,\ldots,n)\) thereafter.
\end{theorem}

The controls are given in \eqref{eq:factorial-control} and
\eqref{eq:near-control}. The positive assertion is
Theorem~\ref{thm:ground}; the negative assertions are
Theorem~\ref{thm:functional-hardness}. Allowing functional inputs
introduces positive ordinal coordinates which carry polynomial-semiring
information. In contrast, with ground free inputs all positive coordinates
are closed, leaving a finite circuit of numerical collapses to compare.

The main new arithmetic mechanism is elementary to state. For
\(p(n)=\lfloor n/2\rfloor\), the controls satisfy
\[
 p(f(n))\equiv p(n)-1\pmod m
\]
on an explicit tail for each \(m\). The greedy collapse evaluator updates
its total norm by \(f\) at every nonfinite event. Every descent beginning
at a head at least \(\omega^2\) must pass through that limit part,
where the remaining event count has a closed formula. The unknown number
of earlier events cancels in the congruence. This yields one uniform
complete-count residue table at every fixed head, however high its
Cantor normal form.

Turning those tables into a decision theorem requires more than
periodicity. Sections~\ref{sec:entry-comparison} and
\ref{sec:arithmetic} prove exact entry comparison, collision identities,
and a finite normalization of signed count rows. The last projection is
ordinary Presburger arithmetic. Closed existential blocks suffice:
failure of a universal comparison is one such block after the finite
assignment circuit is flattened.

The question for the exact control \(G\), or for the exact factorial
control \(G!\), remains separate. Nor is the numerical comparison a
syntactic termination ordering of Bachmann--Howard height: its strict
scalar component at any fixed admissible valuation gives a finite upper
bound on each descending rank. The structural ordering question has a
different proof and is treated in a separate manuscript.

\section{Ordinal norms and effective collapse evaluation}
\label{sec:foundations}

The numerical assignment uses ordinals below $\varepsilon_0$, represented
by hereditary Cantor normal forms.  The ordinal $\varepsilon_0$ is the
least fixed point of $\alpha\mapsto\omega^\alpha$.  We distinguish
ordinary ordinal addition and multiplication from their natural
counterparts $\oplus$ and $\otimes$.  Natural addition collects the
Cantor monomials of its operands; natural multiplication distributes over
those monomials and naturally adds their exponents.  In particular both
natural operations are commutative.  All finite tails displayed as
$\lambda+n$ use ordinary addition.  If $\lambda$ is a limit and $n$ is
finite, then $\lambda\oplus n=\lambda+n$.

\begin{definition}[Additive hereditary norm]\label{def:additive-norm}
For decreasing exponents and positive integer coefficients put
\[
 N(0)=0,\qquad
 N\left(\sum_{i<r}\omega^{\alpha_i}c_i\right)
       =\sum_{i<r}c_i\bigl(1+N(\alpha_i)\bigr).
\]
Thus $N(n)=n$, $N(\omega)=2$ and $N(\omega^2)=3$.
Coefficients are charged with their full multiplicity.
\end{definition}

\begin{lemma}[Effective ordinal operations and norm bounds]
\label{lem:norm-arithmetic}
Every norm ball is finite and effectively enumerable.  Cantor-normal-form
comparison, $\oplus$, $\otimes$, and ordinal exponentiation $2^\alpha$
are effective.  Moreover
\begin{align*}
 N(\alpha\oplus\beta)&=N(\alpha)+N(\beta),\\
 N(\alpha)+N(\beta)-1
 &\le N(\alpha\otimes\beta)\le N(\alpha)N(\beta)
       &&(\alpha,\beta>0),\\
 N(\alpha)&\le2N(2^\alpha),\qquad
 N(2^\alpha)\le2^{N(\alpha)}.
\end{align*}
\end{lemma}
\begin{proof}
An exponent in a nonzero ordinal of norm at most $b$ has norm at most
$b-1$.  Induction on $b$ therefore enumerates the finite set of possible
exponents and the finite lists of coefficients whose total charge is at
most $b$.  Recursive lexicographic comparison decides their order.
The definitions of natural sum and product give finite algorithms.

For the norm bounds, expand coefficients into repeated monomials.
Suppose the two operands have respectively $r,s$ monomials, and let
$A,B$ be the sums of the norms of their exponents.  Their norms are
$r+A$ and $s+B$, while the product has norm $rs+sA+rB$.
Collecting equal monomials preserves these totals.  The upper product
bound follows by expanding $(r+A)(s+B)$; for the lower bound the
difference is
\[
 (r-1)(s-1)+(s-1)A+(r-1)B\ge0.
\]
Norm additivity for natural sum follows in the same way.

Write $\alpha=\omega\xi+n$, $n<\omega$, using ordinary
multiplication.  Then
\[
 2^\alpha=\omega^\xi 2^n,\qquad
 N(2^\alpha)=2^n(1+N(\xi)).
\]
Multiplication on the left by $\omega$ changes the exponent of each
monomial from $\gamma$ to $1+\gamma$.  Its norm stays the same or
increases by one, so
$N(\xi)\le N(\omega\xi)\le2N(\xi)$.
The last two asserted bounds now follow from
$1+a\le2^a$ for $a\in\N$ and
$2N(\xi)+n\le2^{n+1}(1+N(\xi))$.
These are the norm estimates of \cite[Proposition~2.2]{ww}, with the
additive convention made explicit.
\end{proof}

For reference, the finite hierarchy used throughout is
\[
 F_0(n)=2^n,\qquad F_{j+1}(n)=F_j^{\,n+1}(n),\qquad
 G(n)=F_5(n+100).
\]
Iteration is denoted by a superscript.  The proofs below apply to any
computable strictly increasing $f:\N\to\N$ with $f\ge G$.
The controls constructed in Section~\ref{sec:controls} are also primitive
recursive, satisfy the stronger separation estimate
\eqref{eq:growth}, and have an eventual upper bound at a fixed finite
hierarchy level.  These are sufficient conditions within the freedom of
control described in \cite[Section~2.1]{ww}; no assertion that the
resulting assignments equal the assignment for $G$ is intended.

\begin{definition}[Norm-controlled collapse]\label{def:collapse}
For $\alpha<\varepsilon_0$ set
\[
 \psi_f(\alpha)=\max\bigl(\{0\}\cup
   \{\psi_f(\beta)+1:\beta<\alpha,
                    \ N(\beta)\le f(N(\alpha))\}\bigr).
\]
Ordinal recursion defines this function, since each displayed maximum
is finite.
\end{definition}

\begin{lemma}[Basic collapse estimates]\label{lem:collapse-estimates}
If $\beta<\alpha$ and $N(\beta)\le f(N(\alpha))$, then
$\psi_f(\beta)+1\le\psi_f(\alpha)$.  For $n,k\in\N$,
\[
 \psi_f(n)=n,\qquad N(\alpha)\le\psi_f(\alpha),\qquad
 \psi_f(\alpha)+k\le\psi_f(\alpha+k).
\]
If $\alpha\ge\omega$, then $\psi_f(\alpha)>f(N(\alpha))$.
\end{lemma}
\begin{proof}
The first assertion is a term of the defining maximum.  Induction on
$n$ gives the finite values: every predecessor of $n$ is finite and
its norm is below $n\le f(n)$.  If $\alpha\ge\omega$, the finite
ordinal $f(N(\alpha))$ is an admitted predecessor, which proves the
strict estimate and hence the norm bound.  The finite case of the norm
bound has already been proved.  Finally $N(\alpha+k)=N(\alpha)+k$,
so $\alpha+k$ is an admitted predecessor of $\alpha+k+1$.
Induction on $k$ proves the finite-tail inequality.
\end{proof}

In particular a fixed-head collapse is increasing in its finite tail.
The lemma does not assert monotonicity of $\psi_f$ on arbitrary ordinal
arguments.  The norm premise must be verified whenever an ordinal
inequality is passed through the collapse.

\subsection{An exact maximum-predecessor algorithm}

For $\alpha>0$ and $b\in\N$ define
\[
 M(\alpha,b)=\max\{\beta<\alpha:N(\beta)\le b\},\qquad
 L(e,b)=M(\omega^e,b).
\]
The defining set for $M$ contains zero and is finite.  We also put
$M(0,b)=0$ as a convention; it is not a predecessor assertion at zero.

\begin{theorem}[Greedy maximum]\label{thm:maximum-predecessor}
The following simultaneous recursion computes $M$ and $L$ exactly.
Write $\alpha=\omega^e c+r$, where $c>0$ and $r<\omega^e$, and
put $w=1+N(e)$.
\begin{enumerate}[label=(\roman*)]
\item $M(\alpha,0)=0$ and $L(0,b)=L(e,0)=0$.
\item If $r>0$ and $cw\le b$, then
\[
 M(\alpha,b)=\omega^e c+M(r,b-cw).
\]
\item Otherwise set
\[
 d=\begin{cases}
 \min\{c-1,\lfloor b/w\rfloor\},&r=0,\\
 \min\{c,\lfloor b/w\rfloor\},&r>0.
 \end{cases}
\]
Then $M(\alpha,b)=\omega^e d+L(e,b-dw)$, with a zero block
omitted.
\item If $e,b>0$, put $a=M(e,b-1)$, $v=1+N(a)$ and
$(q,s)=\operatorname{divmod}(b,v)$.  Then
\[
 L(e,b)=\omega^a q+L(a,s).
\]
\end{enumerate}
\end{theorem}
\begin{proof}
If an affordable nonzero tail is present, retaining the entire leading
block and decreasing the tail dominates every predecessor with a smaller
leading block.  This proves (ii).  Otherwise the coefficient in (iii)
is the largest affordable coefficient compatible with the strict upper
bound.  Its remaining tail is less than $\omega^e$ and has precisely
the remaining norm budget.

For (iv), a nonzero candidate below $\omega^e$ with budget $b$ has
a leading exponent $a'<e$ satisfying $N(a')\le b-1$.
The largest possible leading exponent is therefore $a=M(e,b-1)$.
Cantor comparison then maximizes its coefficient and the smaller tail
in that order, yielding $q$ and $L(a,s)$.

Termination follows from the lexicographic measure (budget, phase),
with $M$ above $L$ at equal budget.  Clause (ii) lowers the budget;
(iii) either lowers it or passes from $M$ to $L$.  In (iv) the first
recursive call uses $b-1$, while $v\le b$ and $q\ge1$ give $s<b$
for the second.  This proves both effectiveness and exactness.
\end{proof}

\begin{lemma}[Limit saturation and monotonicity]\label{lem:limit-saturation}
For every nonzero limit $\lambda$ and every $b\in\N$,
\[
 N(M(\lambda,b))=b.
\]
For fixed $\lambda$ this maximum is strictly increasing in $b$.
For arbitrary bounds the function $M$ is weakly increasing in each
argument wherever the predecessor interpretation is defined.
\end{lemma}
\begin{proof}
If $\eta=M(\lambda,b)$ has unused budget $d=b-N(\eta)>0$,
then $\eta+d<\lambda$ and $N(\eta+d)=b$, contradicting
maximality.  For $b<c$, the predecessor $M(\lambda,b)+1$ is
below $\lambda$ and has norm at most $c$, proving strict increase.
The last assertion follows by inclusion of the defining candidate sets.
\end{proof}

\begin{theorem}[Single-predecessor collapse recurrence]
\label{thm:single-predecessor}
For every nonzero limit $\lambda$ and $n\in\N$,
\[
 \psi_f(\lambda+n)
   =\psi_f\bigl(M(\lambda,f(N(\lambda)+n))\bigr)+1.
\]
\end{theorem}
\begin{proof}
Induct on $n$.  Put $b=f(N(\lambda)+n)$ and
$\eta=M(\lambda,b)$.  It is an admitted predecessor, of norm $b$,
so it supplies the lower inequality.  Any other admitted $\beta<\lambda$
satisfies $\beta\le\eta$.  If $\beta<\eta$, then
$N(\beta)\le b=N(\eta)\le f(N(\eta))$; controlled strictness
therefore gives $\psi_f(\beta)+1\le\psi_f(\eta)$.

The remaining predecessors are $\lambda+j$, $j<n$.  For such a $j$
put $b_j=f(N(\lambda)+j)$ and $\eta_j=M(\lambda,b_j)$.
Strict growth of $f$ and limit saturation give
$\eta_j<\eta$ and $N(\eta_j)=b_j<b$.
The induction hypothesis and controlled strictness imply
\[
 \psi_f(\lambda+j)=\psi_f(\eta_j)+1\le\psi_f(\eta).
\]
Thus every contribution to the original maximum is at most
$\psi_f(\eta)+1$, which proves the reverse inequality.
\end{proof}

\begin{corollary}[Closed evaluation and its norm orbit]
\label{cor:closed-evaluation}
A closed collapse value is computable by the following deterministic
descent.  At a nonfinite state $\alpha=\lambda+n$, replace it by
$M(\lambda,f(N(\alpha)))$ and add one to an event counter.
On reaching a finite state $Z$, return $Z$ plus that counter.
If the initial norm is $t$ and the number of events is $D$, then
\[
 Z=f^D(t),\qquad \psi_f(\alpha)=Z+D.
\]
For a nonfinite initial state one also has $0<D<Z$.
\end{corollary}
\begin{proof}
Each step strictly lowers the ordinal, so well-foundedness below
$\varepsilon_0$ proves termination.  The preceding theorem preserves
the value after adding the event count, and finite evaluation gives the
last value.  Limit saturation makes the norm of each new state exactly
$f$ of the preceding norm, including the transition to the finite state.
There is no further event count along the finite predecessor chain.
A nonfinite start has $t\ge2$, and $f(n)>n$ gives
$Z\ge t+D>D$.
\end{proof}

\begin{corollary}[Finite coefficient identity]\label{cor:finite-coefficient}
For all $k,n\in\N$,
\[
 \psi_f(\omega k+n)=f^k(2k+n)+k.
\]
\end{corollary}
\begin{proof}
The case $k=0$ is finite evaluation.  If $k>0$ and
$b=f(2k+n)$, the exact maximum is
$M(\omega k,b)=\omega(k-1)+(b-2(k-1))$.
It has total norm $b$.  The single-predecessor recurrence and induction
on $k$ give the formula.
\end{proof}

The evaluator always calls $M$ with the \emph{limit part} of its
current ordinal.  Calling it with $\lambda+n$ instead would retain
finite predecessors that have already been accounted for in the maximum
argument and would not give the stated norm orbit.

\section{Literal canonical assignments and the repaired norm proof}
\label{sec:source}

\subsection{Typed terms and ordinal expression trees}

The ground type is $o$, and arrows associate to the right.  Put
\[
 \lv(o)=0,\qquad
 \lv(\sigma\to\tau)=\max\{\lv(\sigma)+1,\lv(\tau)\}.
\]
Terms of the original system T are variables, applications and
abstractions, together with
\[
 0:o,\quad S:o\to o,\quad
 D^\tau:o\to\tau\to\tau\to\tau,\quad
 R^\tau:o\to(o\to\tau\to\tau)\to\tau\to\tau.
\]
The defining reductions, in arbitrary well-typed contexts, are
\begin{align*}
 (\lambda X.A)B&\longrightarrow A[X:=B],\\
 D^\tau0AB&\longrightarrow A,&
 D^\tau(Sn)AB&\longrightarrow B,\\
 R^\tau0AB&\longrightarrow B,&
 R^\tau(Sn)AB&\longrightarrow A\,n\,(R^\tau nAB).
\end{align*}
Substitution is capture avoiding.  The assignment below reads the
literal typed term; it does not first normalize it or identify
extensionally equal terms.  The annotated recursor constructor used
temporarily in the reduction analysis of \cite{ww} is not a source
constructor here.

For each typed variable $X:\sigma$ introduce a vector
$\mathbf x=(x_0,\ldots,x_{\lv(\sigma)})$ of distinct ordinal
variables.  The ambient expression grammar has leaves $0,1,\omega$
and these component variables, and three binary constructors
\[
 \operatorname{Add}(a,b),\qquad
 \operatorname{ExpProd}(a,b),\qquad
 \operatorname{Psi}(a,b).
\]
Their values are respectively
\[
 a\oplus b,\qquad 2^a\otimes b,\qquad
 \psi_f(\omega\otimes a\oplus b).
\]
An exponential product or collapse is one primitive binary constructor,
not an arbitrary multiplication tree.  We often write $a+b$ for the
literal binary constructor $\operatorname{Add}(a,b)$; this shorthand
does not authorize reassociation.  Interpretation always uses natural
sum.  An expression is $\mathbf x$-free when no component of
$\mathbf x$ occurs in its tree.

A vector of level $m$ records exactly $m+1$ actual coordinate trees.
An out-of-range coordinate is absent and is interpreted as zero in a
comparison.  When vectors are added, two actual coordinates are joined
by one Add node, and a single actual coordinate is retained unchanged.
Thus addition does not insert actual Zero nodes to represent absent
coordinates.  All actual coordinates in the graded classes and in the
source assignment are zero-free and have positive values.

\subsection{The substitution domain and comparisons}

The closed classes $B_i$ are the least simultaneous classes satisfying
the following rules, which are those of \cite[Definition~2.6]{ww}.
\begin{enumerate}[label=(\roman*)]
\item $1\in B_i$ for every $i$, and $\omega\in B_i$ for $i\ge1$.
\item If $a,b\in B_i$, then $\operatorname{Add}(a,b)\in B_i$.
\item If $i\ge1$, $a\in B_{i+1}$ and $b\in B_i$, then
$\operatorname{ExpProd}(a,b)\in B_i$.
\item If $a\in B_1$, $b\in B_0$ and $N(a)\le F_2(b)$, then
$\operatorname{Psi}(a,b)\in B_0$.
\item Every member of $B_0$ is a member of $B_i$ for $i\ge1$.
\end{enumerate}
All expressions in $B_0$ have positive finite values.  A closed vector
$h$ belongs to $B$ when $h_i\in B_i$ at every actual coordinate.
It has \emph{bounded norm} when
\[
 N(h_i)\le N(h_0)=h_0\quad\text{for every }i.
\]

An admissible valuation replaces each occurring variable vector by one
whole vector in $B$ of the same level and of bounded norm.  If even one
component of a vector occurs in either expression being compared, all
components of that vector belong to the domain of the valuation.
Thus the components of a functional variable cannot be varied as
independent unconstrained ordinal parameters.

Scalar relations $a\mle b$, $a\mslt b$ and $a=b$ mean that their
values satisfy $\le$, $<$ and $=$ under every admissible valuation.
The same convention applies to norm inequalities.  For vectors use
the zero padding just specified and set
\[
 a\mle b\ \Longleftrightarrow\ (\forall i)\ a_i\mle b_i,
 \qquad
 a\mslt b\ \Longleftrightarrow\
 a_0\mslt b_0\ \text{ and }\ (\forall i>0)\ a_i\mle b_i.
\]
Equality is universal coordinatewise equality.  These are precisely the
comparison conventions of \cite[Definitions~2.7--2.8]{ww}.
For a ground variable the admissible values are exactly the positive
integers: every such integer is represented by a finite sum of ones,
and every $B_0$ expression has a positive integer value.

To describe the operators, define open classes $C_i$ by the same rules
as $B_i$, with the additional leaf $y_i\in C_i$ whenever
$i\le\lv(Y)$.  The collapse guard is now universal under admissible
valuations.  For a fixed $\mathbf x$, define $C_i^{\mathbf x}$ in
the same way, except that promotion from grade zero to a positive
grade is permitted only for a wholly $\mathbf x$-free expression.
Write $C$ and $C^{\mathbf x}$ for coordinatewise graded vectors.
Their closed instances are exactly $B$.  These classes do not include
a bounded-norm requirement; it is a separate property of a vector.

Induction on this grammar shows that an expression in
$C_i^{\mathbf x}$ contains no $x_j$ with $j<i$.
In particular, if $n=\lv(X)+1$, an expression of grade at least
$n$ is already $\mathbf x$-free.  This stopping property will pay
the norm budget in abstraction.

\subsection{A finite literal parsing policy}

Finite sums and repetitions in \cite[Definitions~2.10 and 2.18]{ww}
need an association convention to determine actual expression trees.
This choice matters when a later abstraction examines those trees.
For example the syntactic size defined below gives
\[
 s_{\mathbf x}(x_0+(1+1))=3,\qquad
 s_{\mathbf x}((x_0+1)+1)=5,
\]
although the two ordinal-valued expressions are equal.

\begin{definition}[Admissible literal policy]\label{def:literal-policy}
A policy $\rho$ is a computable choice of an ordered full binary
addition tree for every nonempty finite-sum macro site and arity.
It uses every operand once, in order and as an unchanged subtree.
Its choice depends on the site and arity, not on operand names,
values, internal trees or free-variable status.  A positive $k$-fold
repeat contains exactly $k$ unchanged copies; a one-fold repeat is
the operand itself.  Existing trees are never flattened or reassociated.
\end{definition}

The sites are positive numerals, positive repetitions, the active
three-operand sum partial, the active exponential partial, positive
coordinate aggregation in abstraction, its scalar repeat, and the
three-operand box tail when its finite contribution is nonzero.
In the exponential partial the first operand is the already formed
two-fold repeat of the head partial; it is not flattened into the
outer sum.  All expressly binary clauses remain binary.

One possible default uses left combs for the unspecified sum and repeat
sites, with box tail $a_0+(b_0+n)$ for $n>0$ and $a_0+b_0$ for
$n=0$.  The theorem allows any supplied policy in Definition~\ref{def:literal-policy}.
It neither requires nor asserts independence of the final assignment
from that policy.  A zero finite contribution in the box is absent;
this convention does not simplify any actual Zero tree.  The source
constant $0$ will be assigned the positive scalar $1$.

\subsection{Application and its domain estimate}

Let $a,b$ have levels $m>n$, and define $c=a\mathbin\square b$
downwards by
\begin{align*}
 c_i&=a_i&& (n<i\le m),\\
 c_i&=\operatorname{ExpProd}(c_{i+1},\operatorname{Add}(a_i,b_i))
                                      &&(1\le i\le n),\\
 c_0&=\begin{cases}
 \operatorname{Psi}(c_1,a_0+(b_0+n)),&n>0,\\
 \operatorname{Psi}(c_1,a_0+b_0),&n=0.
 \end{cases}
\end{align*}
The positive numeral $n$ and any optional association of the three
tail operands are expanded by $\rho$; their value is always
$a_0+b_0+n$.  Restriction $c\upharpoonright_r$ keeps the actual
coordinates $0,\ldots,r$.

\begin{lemma}[Box domain and bounded norm]\label{lem:source-box}
If $a,b\in C$ have bounded norm and $\lv(a)>\lv(b)=n$, then
$a\square b$ belongs to $C$ and has bounded norm.  The same statement
holds with $C^{\mathbf x}$ in place of $C$.
\end{lemma}
\begin{proof}
Fix an admissible valuation and put $u=a_0+b_0+n\ge n+2$.
All positive input-coordinate norms are at most $u$.
The product and exponential bounds in Lemma~\ref{lem:norm-arithmetic}
show that, at one downward step, a bound $x\ge u$ is replaced by
at most $u2^x$.  Since $u2^x\le2^{2^x}$ for $x\ge u\ge2$,
after at most $n$ steps every positive output norm is at most
$F_0^{2n}(u)$.  Directly from the iteration definitions,
\[
 F_0^{2n}(u)\le F_0^{2(u+1)}(u)\le F_1^2(u)\le F_2(u).
\]
For the middle inequality, expand both successive evaluations of
$F_1(v)=F_0^{v+1}(v)$; each contributes at least $u+1$
iterations beginning at a value at least $u$.

Consequently $N(c_1)\le F_2(u)$, which is the actual grammar guard
for $c_0=\operatorname{Psi}(c_1,u)$.  The positive coordinates
have their required grades by downward induction.  Put
$\alpha=\omega\otimes c_1\oplus u$.  It is nonfinite and
$N(\alpha)\ge u$, whence
\[
 c_0=\psi_f(\alpha)>f(N(\alpha))\ge G(u)\ge F_2(u).
\]
This bounds the norms of every positive output by its scalar output.
All estimates hold at every admissible valuation.  No promotion clause
is added in this construction, so the same reasoning proves the
$C^{\mathbf x}$ assertion.
\end{proof}

\subsection{The partials and the unchanged abstraction budget}

Fix $\mathbf x$ and put $n=\lv(X)+1$.  The syntactic size tests
whole-freeness first:
\[
 s_{\mathbf x}(h)=
 \begin{cases}
 1,&h\text{ is wholly }\mathbf x\text{-free, or }h=x_i,\\
 s_{\mathbf x}(a)+s_{\mathbf x}(b)+1,
    &h=\operatorname{Add}(a,b)\text{ or }\operatorname{Psi}(a,b),\\
 2s_{\mathbf x}(a)+s_{\mathbf x}(b)+1,
    &h=\operatorname{ExpProd}(a,b).
 \end{cases}
\]
The last two clauses apply only if the first does not.
For $h\in C_i^{\mathbf x}$, the vector $\partial_i^{\mathbf x}h$
has level $n$.  If $h$ is wholly $\mathbf x$-free, set
\[
 (\partial_i^{\mathbf x}h)_j=
 \begin{cases}h+1,&j=i\le n,\\1,&j\ne i\text{ or }i>n.
 \end{cases}
\]
For active expressions the clauses are
\begin{align*}
 \partial_i^{\mathbf x}x_i&=\mathbf1,\\
 \partial_i^{\mathbf x}\operatorname{Add}(a,b)
   &=\partial_i^{\mathbf x}a+\partial_i^{\mathbf x}b+\mathbf1,\\
 \partial_i^{\mathbf x}\operatorname{ExpProd}(a,b)
   &=2\partial_{i+1}^{\mathbf x}a+\partial_i^{\mathbf x}b+\mathbf1
       \qquad(i\ge1).
\end{align*}
All displayed vectors $\mathbf1$ have level $n$.  Multiples and
three-operand sums use the policy, including the nested two-fold head
repeat in the last clause.  For an active scalar collapse, put
\begin{align*}
 (\partial_0^{\mathbf x}\operatorname{Psi}(a,b))_0
   &=\operatorname{Psi}(a\{\mathbf x:=\mathbf1\},
                              (\partial_0^{\mathbf x}b)_0),\\
 (\partial_0^{\mathbf x}\operatorname{Psi}(a,b))_j
   &=(\partial_1^{\mathbf x}a)_j+(\partial_0^{\mathbf x}b)_j
                                      \qquad(1\le j\le n).
\end{align*}

For $h\in C^{\mathbf x}$ of level $m$, let
\[
 S_{\mathbf x}(h)=2^n\sum_{i=0}^m s_{\mathbf x}(h_i).
\]
This is a natural number computed from the fixed input trees.
The vector $\delta^{\mathbf x}h$ has level $\max\{n,m\}$ and
coordinates
\begin{equation}\label{eq:source-delta}
 (\delta^{\mathbf x}h)_j=
 \begin{cases}
 S_{\mathbf x}(h)\,(\partial_0^{\mathbf x}h_0)_0,&j=0,\\
 \displaystyle\sum_{i=0}^m(\partial_i^{\mathbf x}h_i)_j,&0<j\le n,\\
 h_j,&n<j\le m.
 \end{cases}
\end{equation}
The scalar multiplication is the literal positive repetition at its
designated policy site.  The natural number $S_{\mathbf x}(h)$ is
not recomputed after reassociating any input tree.

\begin{lemma}[Finite definition and substitution covariance]
\label{lem:source-finiteness}
For a supplied policy these definitions produce unique finite trees
effectively.  The partials and $\delta^{\mathbf x}$ commute with
substitutions for other variable vectors by $\mathbf x$-free vectors,
whenever the displayed operations have their graded domains.
\end{lemma}
\begin{proof}
Each partial is defined by structural recursion on its input tree.
For an explicit finite bound, let $|h|$ be its node count.  Give a leaf
weight one, Add and Psi weight $w(a)+w(b)+1$, and ExpProd weight
$2w(a)+w(b)+1$.  Induction gives
$s_{\mathbf x}(h)\le w(h)\le2^{|h|}-1$.
Every coordinate of a partial has at most $6\cdot3^{|h|}$ nodes:
the active Add, ExpProd, scalar Psi and positive Psi clauses have
bounds respectively
\[
 E(a)+E(b)+3,\quad2E(a)+E(b)+4,\quad
 |a|+E(b)+1,\quad E(a)+E(b)+1,
\]
where $E(a)=6\cdot3^{|a|}$, and the whole-free diagonal has
$|h|+2$ nodes.  Each fits the asserted recursive bound.
The finite sums and repeats in \eqref{eq:source-delta} therefore
have computable sizes.  This construction never has to decide a
universal comparison in order to construct a tree.

A substitution by $\mathbf x$-free vectors preserves whether a
subtree contains a component of $\mathbf x$.  Whole-free subtrees
still contribute size one, while the active constructor pattern is
unchanged.  Hence it preserves $s_{\mathbf x}$ and
$S_{\mathbf x}$.  Induction on each partial clause proves literal
commutation: the policy depends only on the unchanged macro site and
arity.  The finite sums in $\delta^{\mathbf x}$ then commute too.
\end{proof}

\begin{lemma}[Weak scalar estimate and the graded domain]
\label{lem:weak-scalar}
For $h\in C_0^{\mathbf x}$,
\[
 h\{\mathbf x:=\mathbf1\}\mle
             (\partial_0^{\mathbf x}h)_0.
\]
Every $\partial_i^{\mathbf x}h$, for $h\in C_i^{\mathbf x}$,
is an $\mathbf x$-free $C$-vector of level $n$.
For $h\in C^{\mathbf x}$, the vector $\delta^{\mathbf x}h$
is also valid and $\mathbf x$-free.
\end{lemma}
\begin{proof}
Induct simultaneously on the grammar and the partial recursion.
A whole-free scalar has diagonal $h+1$; the tracked scalar variable
has diagonal one; an active Add has the two scalar partials plus one.
For an active Psi the two heads after replacing $\mathbf x$ by ones
are identical, and the induction hypothesis compares their finite tails.
Lemma~\ref{lem:collapse-estimates} supplies the required weak
finite-tail monotonicity.

That new Psi has its grammar guard: the old guard and the weak
induction hypothesis give
\[
 N(a\{\mathbf x:=\mathbf1\})
 \mle F_2(b\{\mathbf x:=\mathbf1\})
 \mle F_2((\partial_0^{\mathbf x}b)_0).
\]
The head has the correct positive grade.  Whole-free diagonal outputs
are $h+1$ at their original grade; all off-diagonal and tracked-variable
outputs are one.  The other positive outputs only add correctly graded
partials, with the displayed fixed multiplicities.  Every output is
$\mathbf x$-free.  Sums and positive scalar repetitions preserve
these properties.  Above grade $n$, the retained $h_j$ is already
$\mathbf x$-free by the stopping property of the grammar.
\end{proof}

\begin{remark}[Two strict inequalities that are unnecessary]
The strict form printed in \cite[Lemma~2.19]{ww} fails at $h=x_0$:
both scalar sides are one.  Its appendix estimate (A.5) also has an
equality case: take $\lv(X)=0$, $h=(x_0,1)$, the coordinate
$i=1$, and its free subtree $1$.  The asserted strict estimate becomes
$1<1$.  The weak scalar lemma just proved is sufficient for the
collapse guard.  The next two results replace the strict free-subtree
estimate while preserving the multiplier in \eqref{eq:source-delta}.
These observations do not refute the separate beta estimate involving
both $\delta$ and $\square$ in \cite[Lemma~2.22]{ww}.
\end{remark}

\begin{lemma}[First free positive subtree]
\label{lem:first-free-subtree}
Let $h\in C^{\mathbf x}$ have bounded norm.  Fix an admissible
valuation for its other names and put
$B=h_0\{\mathbf x:=\mathbf1\}$ at that valuation.
In a partial recursion beginning at any actual $h_i$, every first
wholly $\mathbf x$-free subtree $a$ encountered at a positive
grade satisfies
\[
 N(a)\le2^{n-1}B.
\]
\end{lemma}
\begin{proof}
Measure norms after replacing $\mathbf x$ by ones.  The free
subtree is unchanged, and every initial coordinate then has norm at
most $B$.  A step into an Add child cannot increase norm.  A step
into an exponential tail cannot increase it either, by positivity and
the lower product bound.  A step into an exponential head costs at
most a factor two, since
\[
 N(a)\le2N(2^a)\le2N(2^a\otimes b).
\]
It also raises the grade by one.  At grade $n$ a subtree is already
wholly $\mathbf x$-free and recursion stops.  A path beginning
at positive grade thus makes at most $n-1$ such head steps.

A scalar path enters a positive grade only through a collapse head.
This step costs no factor two:
\[
 N(a)\le N(\omega\otimes a\oplus b)
       \le\psi_f(\omega\otimes a\oplus b).
\]
Its scalar tail is bounded by the same collapse value, so scalar
tail steps also cost no factor.  The head enters grade one, leaving
at most $n-1$ active exponential-head steps.  A scalar promoted to
positive grade is already wholly free, so it cannot begin another
active scalar descent.  Multiplication of the possible factors two
proves the bound.
\end{proof}

\begin{theorem}[Bounded norm with the original multiplier]
\label{thm:delta-bounded-norm}
If $h\in C^{\mathbf x}$ has bounded norm, then
$\delta^{\mathbf x}h$ has bounded norm, with exactly the multiplier
$S_{\mathbf x}(h)=2^n\sum_i s_{\mathbf x}(h_i)$ in
\eqref{eq:source-delta}, for every admissible policy.
\end{theorem}
\begin{proof}
Fix the other names as in Lemma~\ref{lem:first-free-subtree} and put
$H=(\partial_0^{\mathbf x}h_0)_0$ and $v=2^nH$.
The weak scalar estimate gives $1\le B\le H$.
Whenever a free positive-grade subtree can be copied into a positive
diagonal output, the preceding lemma gives
\[
 N(a)+1\le2^{n-1}B+1\le2^nH=v.
\]
A free \emph{scalar} subtree is not copied into a positive output;
those outputs are one.  This distinction supplies the extra unit in
the whole-free diagonal clause without a strict free-subtree bound.

Induct through the actual partial recursion, using those bounds at
each first free leaf.  For $0<j\le n$ the result is
\[
 N((\partial_i^{\mathbf x}h_i)_j)
           \le s_{\mathbf x}(h_i)v.
\]
Tracked-variable outputs have norm one.  An active Add contributes
the two child bounds and one, bounded by
$(s_{\mathbf x}(a)+s_{\mathbf x}(b)+1)v$.
An active ExpProd partial has two copies of its head output, one
tail output and one, bounded by
$(2s_{\mathbf x}(a)+s_{\mathbf x}(b)+1)v$.
A positive Psi output only adds its head and tail partials, so its
size allowance also suffices.  Norm additivity makes every allowed
association of these sums satisfy the same estimate.

Summing over $i$ gives
\[
 N((\delta^{\mathbf x}h)_j)
 \le\left(\sum_i s_{\mathbf x}(h_i)\right)2^nH
   =(\delta^{\mathbf x}h)_0\qquad(0<j\le n).
\]
For $j>n$ the retained $h_j$ is $\mathbf x$-free and its norm
is at most $B\le H\le(\delta^{\mathbf x}h)_0$.
Since the valuation was arbitrary, this proves bounded norm.
\end{proof}

\subsection{The canonical assignment and its arithmetic image}

The following is the zero-history instance of
\cite[Definition~3.1]{ww}, with literal expansions fixed by $\rho$:
\begin{align*}
 \can{X:\sigma}&=\mathbf x,\qquad
 \can0=(1),\qquad \can S=(1,1),\\
 \can{D^\tau}&=(1,\ldots,1)
                &&\text{of level }\lv(\tau)+1,\\
 \can{R^\tau}&=(2,1,\ldots,1,\omega)
                &&\text{of level }\lv(\tau)+2,\\
 \can{AB}&=(\can A\square\can B)\upharpoonright_{\lv(\tau)}
                &&(A:\sigma\to\tau, B:\sigma),\\
 \can{\lambda X.G}&=\delta^{\mathbf x}\can G+
                        \can G\{\mathbf x:=\mathbf1\}.
\end{align*}
The numeral $2$ is the literal $1+1$.  The vector of ones in a
substitution has the level of the variable being replaced.  When useful
we write $\can A_{f,\rho}$ to display the two fixed choices.

\begin{proposition}[Well-defined canonical source assignment]
\label{prop:canonical-domain}
Every literal original T term has an effectively constructed canonical
vector of the level of its type.  It belongs to $C$ and has bounded
norm.  A body to which $\delta^{\mathbf x}$ is applied belongs
to $C^{\mathbf x}$.  All these assertions allow arbitrary finite
internal and bound-variable types and arbitrary recursor instances.
\end{proposition}
\begin{proof}
The construction is structural recursion on the source term and the
finite recursions of Lemma~\ref{lem:source-finiteness}.  If
$A:\sigma\to\tau$, then
$\lv(A)\ge\lv(\sigma)+1>\lv(\sigma)$, which is exactly the
level premise for the box.  Restriction leaves level $\lv(\tau)$;
the abstraction has level $\max\{\lv(X)+1,\lv(G)\}$ as required.

For its domain, keep the following strengthened invariant: a positive
coordinate is generated from $1,\omega$, appropriate positive-index
component variables, Add and graded ExpProd, without any scalar
promotion or Psi.  The constants and variables have this property.
The positive box clauses preserve it.  In a scalar partial a whole-free
subtree contributes only ones to positive outputs, an active variable
contributes ones, and Add joins partials; an active Psi joins the
positive partials of its positive head and scalar tail.  In a partial
of a positive input, whole-free copies, variables, Add and ExpProd
also preserve the invariant.  The sums and retained coordinates in
abstraction do so as well.  This proves the invariant simultaneously
with source construction.  In particular every positive coordinate
has the grammar required for $C_i^{\mathbf x}$ for each tracked
binder, and its scalar collapse guards follow from
Lemmas~\ref{lem:source-box} and \ref{lem:weak-scalar}.

Variables have bounded norm by the valuation domain.  The constants
do too; for the recursor its top norm is $N(\omega)=2$, equal to
its scalar.  The box preserves bounded norm by
Lemma~\ref{lem:source-box}, as does restriction.  At an abstraction
the first summand has bounded norm by
Theorem~\ref{thm:delta-bounded-norm}; the second does because the
all-ones vector is admissible.  Coordinatewise natural addition
preserves the bound by norm additivity.  This completes the induction.
\end{proof}

\begin{theorem}[Fixed-head arithmetic image of actual sources]
\label{thm:generated-source}
Let $A$ be an original T term whose free variables all have ground
type.  For every positive coordinate, $\can A_i$ is a closed
effectively computable ordinal.  Its scalar is an effectively obtained
finite circuit in the ground component inputs, positive integer
constants, addition, fixed positive integer multiples, and functions
\[
 h_\lambda(u)=\psi_f(\lambda+u)
\]
with fixed nonzero limit heads $\lambda<\varepsilon_0$.
Repeated occurrences of one source variable give one shared input;
shared subcircuits may also be retained.  The theorem imposes no bound
on internal types, bound-variable types, result type, recursors, or
finite nesting, and is uniform in the supplied computable policy.
\end{theorem}
\begin{proof}
The positive-coordinate invariant in the preceding proof contains
only positive-index components of variables that remain free.
Every bound component is removed by the partials or by the simultaneous
all-ones substitution at its abstraction.  Ground free variables have
no positive-index components.  All positive output coordinates are
therefore closed, and their values are computable by the finite ordinal
operations of Lemma~\ref{lem:norm-arithmetic}.

The same induction shows that an actual scalar coordinate is formed
only from scalar component variables, positive constants, Add,
fixed repeats and Psi.  Every head of a scalar Psi is a positive
expression in the invariant just proved.  After the ground-free
restriction its value is a fixed closed ordinal $c>0$.
Set $\lambda=\omega\otimes c$.  Every monomial of this ordinal
has positive exponent, so $\lambda$ is a nonzero limit, and its
finite scalar tail $u$ gives
$\operatorname{Psi}(c,u)=\psi_f(\lambda+u)$.
The syntactic multiplier $S_{\mathbf x}(h)$ is a fixed integer
computed before valuations, so it adds no new variable or quantifier.
Translate this finite tree recursively, retaining names and shared
wires.  This is the required circuit.

Its head is the one produced by the literal source assignment.
The argument does not replace a generated head by an arbitrary ordinal
primitive or require a source term for every auxiliary head later used
by arithmetic elimination.
\end{proof}

The zero-history assignment is the object of comparison here.  The
normalization theorem in \cite{ww} permits history-dependent assignment
derivations at abstractions in order to follow reductions.  No claim
that this single fixed canonical map decreases at every reduction is
needed for the comparison theorem.

\section{A control with a one-step modular potential}\label{sec:controls}

The freedom to choose the numerical control is decisive.  We use the
finite hierarchy
\[
 F_0(n)=2^n,\qquad F_{j+1}(n)=F_j^{\,n+1}(n),\qquad
 G(n)=F_5(n+100).
\]
The superscript on a function denotes iteration.  Put
\(p(n)=\lfloor n/2\rfloor\) and
\(L(n)=\operatorname{lcm}(1,\ldots,n)\), with \(L(0)=1\).
For a prescribed integer \(b\ge0\), define
\begin{align}
 f_*(n)&=G(n)!+2p(n)-2,\label{eq:factorial-control}\\
 f_b(n)&=\begin{cases}
 G(n),&n\le b,\\
 G(n)+((2p(n)-2-G(n))\bmod 2L(n)),&n>b.
 \end{cases}\label{eq:near-control}
\end{align}
All remainders are least nonnegative remainders.  The first definition
makes the divisibility transparent.  The second preserves an arbitrary
finite prefix of the published example and satisfies
\begin{equation}\label{eq:control-distance}
 0\le f_b(n)-G(n)<2L(n)\le2n!\quad(n>b).
\end{equation}
Below, \(f\) is either of these controls, and
\(B_f=0\) for \(f_*\), \(B_f=b+1\) for \(f_b\).

\begin{lemma}[Growth and modular decrement]\label{lem:control}
The functions \(f_*\) and \(f_b\) are even-valued primitive recursive
functions, majorize \(G\), and satisfy
\begin{equation}\label{eq:growth}
 f(n+1)\ge2^{2^{f(n)}},\qquad f(n)<G(n+1).
\end{equation}
They are eventually bounded by one fixed finite level of the hierarchy.
For every \(m\ge1\),
\begin{equation}\label{eq:potential}
 p(f(n))\equiv p(n)-1\pmod m
 \qquad(n\ge\max\{m,B_f\}).
\end{equation}
\end{lemma}
\begin{proof}
Factorial, least common multiple, and bounded remainder are primitive
recursive.  Every value of \(G\) is even.  The defining congruence for
\(f_b\) is even, so \(f_b\) is even too.  Write \(A=G(n)\).
Monotonicity of the iterates defining \(F_5\), including the extra iterate
at the larger argument, gives
\[
 G(n+1)\ge F_4(G(n)+1).
\]
For \(A\ge16\),
\[
 2A!\le2A^A<2^{A^2}<2^{2^A}.
\]
The fixed shift by 100 ensures this range at every input.  Also
\(G(n)\ge2n!\) and \(G(n)\ge2n+2\).  Consequently
\[
 G(n)\le f_b(n)<2G(n),\qquad
 G(n)\le f_*(n)\le2G(n)!.
\]
The latter lower bound holds also at \(n=0,1\), where subtracting two
from \(G(n)!\) still leaves a number above \(G(n)\).
The function \(F_4(A+1)\) dominates four successive exponentiations of
two beginning at \(A\).  The preceding estimates therefore imply
\eqref{eq:growth}, including the transition from the preserved prefix
to the corrected tail in \(f_b\).  For all sufficiently large \(n\),
\(G(n+1)=F_5(n+101)\le F_5(F_5(n))\le F_6(n)\), which is the
required eventual finite-level bound.

For \(f_*\) one has
\[
 p(f_*(n))=G(n)!/2+p(n)-1.
\]
If \(m\le n\), the integer \(2m\) is at most \(G(n)\), and hence
\(m\mid G(n)!/2\).  For \(f_b\), its definition gives
\(f_b(n)\equiv2p(n)-2\pmod{2L(n)}\) on the corrected tail.
Divide this congruence by two and use \(m\mid L(n)\).  This proves
\eqref{eq:potential}.
\end{proof}

The estimate \eqref{eq:growth} supplies the norm and separation margins
used in the canonical assignment.  Equation \eqref{eq:potential} has a
different purpose: it records the number of descent events without
computing that number.  Since the norm orbit is strictly increasing,
ordinary induction gives
\begin{equation}\label{eq:orbit-potential}
 p(f^i(t))\equiv p(t)-i\pmod m
 \quad\text{for every }i\ge0,
 \quad t\ge\max\{m,B_f\}.
\end{equation}
There is no bound on \(i\) in this identity.

\section{Complete descent counts and entry residues}\label{sec:residues}

For a nonzero limit \(\lambda<\varepsilon_0\), define
\[
 J_\lambda(t)=\psi_f(\lambda+t-N(\lambda)),
 \qquad t\ge N(\lambda).
\]
The greedy evaluation starts at \(\lambda+t-N(\lambda)\).  At a
nonfinite state \(h+k\), where \(h\) is its nonzero limit part, it
moves to \(M(h,f(N(h+k)))\).  Let \(D_\lambda(t)\) be the number
of such moves and \(Z_\lambda(t)\) the first finite state.  Thus the
finite predecessor chain after that state is not included in \(D\).
The maximal-predecessor and single-path lemmas give
\begin{equation}\label{eq:terminal-decomposition}
 u_i=f^i(t),\qquad Z_\lambda(t)=f^{D_\lambda(t)}(t),\qquad
 J_\lambda(t)=Z_\lambda(t)+D_\lambda(t).
\end{equation}
The use of the limit part in the predecessor call is essential.

\begin{theorem}[Uniform complete-count residues]\label{thm:counts}
For every \(m\ge1\) and every nonzero limit \(\lambda\), put
\(T=\max\{N(\lambda),m,B_f\}\).  For all \(t\ge T\),
\begin{align*}
 \lambda\ge\omega^2:\quad&
 D_\lambda(t)\equiv p(t)\pmod m,\quad
 Z_\lambda(t)\equiv0\pmod{2m},\quad
 J_\lambda(t)\equiv p(t)\pmod m;\\
 \lambda=\omega k:\quad&
 D_\lambda(t)=k,\quad
 Z_\lambda(t)\equiv2p(t)-2k\pmod{2m},\quad
 J_\lambda(t)\equiv2p(t)-k\pmod m.
\end{align*}
In particular, the first line gives a period \(2m\) and the explicit
table \(r\mapsto\lfloor r/2\rfloor\bmod m\) for every head
\(\lambda\ge\omega^2\).
\end{theorem}
\begin{proof}
Suppose \(\lambda\ge\omega^2\).  As long as the current limit part
is above \(\omega^2\), the ordinal \(\omega^2\), of norm three,
is an affordable predecessor.  The next limit part cannot fall below it.
Strict descent of limit parts therefore forces a visit to
\(\omega^2\).  Suppose this visit occurs at step \(i\), with norm
\(v=f^i(t)\).  Evenness of the next budget gives the exact predecessor
\[
 M(\omega^2,f(v))=\omega\cdot f(v)/2.
\]
An \(\omega\)-coefficient decreases by exactly one at each subsequent
nonfinite move, regardless of the finite tail.  Hence
\begin{equation}\label{eq:funnel-count}
 D_\lambda(t)=i+1+f(v)/2=i+1+p(f(v)).
\end{equation}
Equation \eqref{eq:orbit-potential} cancels the unknown \(i\), giving
\(D_\lambda(t)\equiv p(t)\).  Apply that equation once more at the
terminal event.  Since the terminal state is even,
\(p(Z_\lambda(t))\equiv0\pmod m\) implies
\(Z_\lambda(t)\equiv0\pmod{2m}\).  Now use
\eqref{eq:terminal-decomposition}.

For \(\lambda=\omega k\), exactly \(k\) nonfinite moves are made.
The same orbit identity and terminal decomposition give the second line.
These two cases exhaust the nonzero limits below \(\varepsilon_0\).
\end{proof}

The theorem uses one final passage through \(\omega^2\), rather than
an induction computing the count at each higher ordinal.  The actual
count can be enormous; its residue depends only on the original input.

For \(\lambda>\nu\), with \(\lambda\) a nonzero limit and
\(\nu\) a limit or zero, write \(E_{\lambda,\nu}(t)\) for the norm
at the first entry into \([\nu,\nu+\omega)\), and
\(\tau_{\lambda,\nu}(t)\) for its event count.  These functions
are defined only when that entry occurs.  On an affordable tail,
\begin{equation}\label{eq:entry-orbit}
 E_{\lambda,\nu}(t)=f^{\tau_{\lambda,\nu}(t)}(t).
\end{equation}
The case \(\nu=0\) is the terminal pair \((Z_\lambda,D_\lambda)\).
An identity interval is a separate convention with \(E(t)=t\) and
\(\tau(t)=0\).

\begin{proposition}[Entry tables]\label{prop:entry-residues}
Put \(c_\nu=\lceil N(\nu)/2\rceil\).  On a computable tail,
an interval of the form \(\lambda=\nu+\omega k\), \(k\ge1\),
has \(\tau(t)=k\) and \(E(t)=f^k(t)\).  Every other proper interval
has
\begin{equation}\label{eq:long-entry-residues}
 \tau(t)\equiv p(t)-c_\nu\pmod m,\qquad
 E(t)\equiv2c_\nu\pmod{2m}.
\end{equation}
A sufficient tail for these formulas is
\(t\ge\max\{N(\lambda),N(\nu)+3,m,B_f\}\).
\end{proposition}
\begin{proof}
Write \(\nu=P+\omega j\), where every exponent in \(P\), if any,
is at least two.  A proper interval is either the displayed finite
\(\omega\)-tail or has \(\lambda\ge P+\omega^2\).
In the latter case require
\[
 f(t)\ge N(P)+\max\{3,2j\}.
\]
This covers the norm bounds of the intermediate limit \(P+\omega^2\)
and the true target.
The stated tail implies this condition.  If the funnel is reached after
\(i\) events at norm \(v=f^i(t)\), its next \(\omega\)-coefficient
is \(f(v)/2-\lceil N(P)/2\rceil\).  Therefore
\[
 \tau(t)=i+1+p(f(v))-\lceil N(P)/2\rceil-j.
\]
Since \(N(\nu)=N(P)+2j\), the potential identity gives the first
congruence.  Applying it at the entry event gives the second.
For a finite \(\omega\)-tail the coefficient decreases by one at every
event, giving the exact formulas directly.
\end{proof}

For use in arithmetic elimination, the complete tables are as follows.
Here \(a\ge0\), and every displayed value residue has modulus \(2m\)
except in the last row, whose modulus is \(m\).
\begin{center}\small
\begin{tabular}{@{}lll@{}}\toprule
Field & \(k\)-event interval & Other proper interval\\\midrule
\(\tau(t)\pmod m\) & \(k\) & \(p(t)-c_\nu\)\\
\(E(t)\pmod{2m}\) & \(2p(t)-2k\) & \(2c_\nu\)\\
\(f^a(E(t))\pmod{2m}\) & \(2p(t)-2(k+a)\) & \(2c_\nu-2a\)\\
\(f(E(t))/2\pmod m\) & \(p(t)-k-1\) & \(c_\nu-1\)\\\bottomrule
\end{tabular}
\end{center}
For a bare iterate \(f^a(t)\) with \(a\ge1\), the residue modulo
\(2m\) is \(2p(t)-2a\).  The identity field remains \(t\), not
\(2p(t)\); in particular its odd bit must be retained.

The finite prefix cannot be discarded.  If \(b\ge3\), then at
\(t=3\) the preserved control gives
\(D_{\omega^2}(3)=G(3)/2+1\equiv0\pmod3\), whereas
\(p(3)=1\).  The value \(G(3)\) is a power of four, which proves
the displayed congruence without evaluating it.  Small inputs will be
handled by the exact evaluator, not by extending the tail tables.

\section{Effective comparison of entry maps}\label{sec:entry-comparison}

Residues alone do not decide inequalities between sums of collapse values.
The order argument also needs exact comparisons and collision conditions
for entry norms. This section supplies them for the chosen controls.
All ordinal additions in interval endpoints and relative states are ordinary
ordinal additions. Natural sum and product remain explicitly marked.

For a proper interval \(P=(\lambda,\mu)\), let \(\delta_P\) be the
unique nonzero limit with \(\mu+\delta_P=\lambda\). Put
\(E_P=E_{\lambda,\mu}\) and \(\tau_P=\tau_{\lambda,\mu}\).
An effective sufficient domain for entry is
\begin{equation}\label{eq:affordable-tail}
 t\ge\max\{2,N(\lambda),N(\mu)\}.
\end{equation}
Indeed, before a hit the target is an affordable predecessor, so the
descent cannot skip below it; strict ordinal descent forces a hit.
This argument also covers \(\mu=0\). A smaller genuine input may
skip a positive target and is not assigned an invented entry value.

\subsection{Relative descent and the two comparison rules}

For \(y>0\), let \(e(y)\) be its leading Cantor exponent. Define
\[
 s_\mu(e)=N(\mu_{\ge e}),\qquad s_\mu(e(0))=N(\mu).
\]
The subscript means restriction to Cantor terms of exponent at least
\(e\); the convention at zero is separate. Directly from Cantor
normal form,
\begin{equation}\label{eq:relative-norm}
 N(\mu+y)=N(y)+s_\mu(e(y)),\qquad 0\le s_\mu(e(y))\le N(\mu).
\end{equation}
The terms of \(\mu\) below the leading exponent of \(y\) are absorbed;
the contributions at equal exponents add. Before entry, write the actual
state uniquely as \(\mu+y_i\). With \(b_i=f^i(t)\), its exact recursion is
\begin{equation}\label{eq:relative-descent}
 y_0=\delta_P+t-N(\lambda),\qquad
 y_{i+1}=\max\{y<\operatorname{lim}(y_i):
          N(y)+s_\mu(e(y))\le b_{i+1}\}.
\end{equation}
Here \(\operatorname{lim}\) removes the finite tail. The first finite
\(y_i\) is the entry event, and \(N(y_i)\le b_i\).

Consider two recursions of this form with profiles \(s,s'\) bounded by
\(S\), budgets \(b,b'\), and limit parts \(U_i,V_i\).
Two distinct budget comparisons are needed.

\begin{lemma}[Propagation]\label{lem:propagation}
If \(V_{i+k}>U_i=\rho>0\) and
\(b'_{i+k+1}\ge b_i+S\), then
\(V_{i+k+1}\ge\rho>U_{i+1}\).
If instead \(V_{i+k}=U_i\) and
\(b'_{i+k+1}\ge b_{i+1}+S\), then
\(V_{i+k+1}\ge U_{i+1}\).
For equal next budgets the latter conclusion also follows from
\(s'\le s\) on all eligible candidates.
\end{lemma}
\begin{proof}
In the strict case, \(\rho\) is below the upper limit part and
\(N(\rho)+s'(e(\rho))\le N(y_i)+S\le b_i+S\).
It is therefore an upper candidate. The lower next state is strictly
below \(\rho\). In the equal-limit case every lower candidate costs
at most \(b_{i+1}+S\) under the upper profile. Candidate inclusion
also proves the equal-budget alternative.
\end{proof}

If the strict rule persists until the lower process stops at \(n\), the
upper process stops no earlier than \(n+k+1\): at the last lower event
its upper limit part is still positive. The weak rule gives the bound
\(n+k\). Notice that the strict rule compares the next upper budget
with the current lower budget. Using the equal-limit budget condition
for two identical budget orbits would lose the required slack.

\begin{lemma}[Gap and input comparison]\label{lem:gap-input}
For two fixed intervals \(P,Q\), put
\(S=\max\{1,N(\mu_P),N(\mu_Q)\}\). On effective tails:
\begin{enumerate}[label=(\roman*)]
\item If \(\delta_P>\delta_Q\), then
\(\tau_P(t)>\tau_Q(t)\) and \(E_P(t)>E_Q(t)\).
\item If \(\delta_P\ge\delta_Q\) and \(t>u\), then
\(\tau_P(t)\ge\tau_Q(u)\) and \(E_P(t)>E_Q(u)\).
\item For fixed \(k\ge1\), if
\(\delta_P\ge\delta_Q+\omega k\) and \(u+S\le f^k(t)\), then
\(\tau_P(t)\ge\tau_Q(u)+k\) and \(E_P(t)>E_Q(u)\).
\end{enumerate}
\end{lemma}
\begin{proof}
Take initially
\(T\ge\max\{2,S,N(\lambda_P),N(\lambda_Q)\}\).
The estimate \(f(v)\ge v+S\) for \(v\ge T\) propagates the strict
initial gap in (i). For (ii), adjacent growth gives
\(f(t)\ge f(u)+S\); monotonicity and noncontraction preserve this
gap at every later iterate. The weak rule propagates the initial weak
limit comparison. The larger starting input then makes the terminal
norm strictly larger.

For (iii) enlarge the tail by \(2k\). During its first \(k\) events
the upper process retains the successive reserves
\(\delta_Q+\omega(k-j)\). At event \(j<k\) the next reserve costs
at most \(N(\delta_Q)+2(k-j-1)+S\), below the next norm budget.
At lag \(k\), compare with the lower process using the weak rule and
\(f^k(t)\ge u+S\). This yields the count bound. Since \(S\ge1\),
the terminal inequality is strict. The positive unit is necessary when
both targets are zero: with \(P=(\omega\cdot2,0)\), \(Q=(\omega,0)\),
and \(u=f(t)\), the terminal norms coincide.
\end{proof}

Distinct limit gaps differ by at least an \(\omega\)-tail. Thus fixed
input translations do not change their eventual order, by (iii) with
\(k=1\). For equal gaps, unequal fixed input translations are ordered
by (ii). It remains to compare equal gaps at one common input.

\subsection{Equal gaps: a finite phase calculation}

Suppose \(\delta_P=\delta_Q=\delta\). Set
\[
 R=20+N(\delta)+N(\mu_P)+N(\mu_Q),\qquad
 M=\operatorname{lcm}(1,\ldots,R+1).
\]
Choose the input tail above \(100(R+2)^4\), \(M\), and \(B_f\),
and above the finitely many affordability bounds used below. In
particular every next tail budget exceeds the current prefix cost and
all fixed shifts by more than \((R+3)^3\).
All fixed exponents and breakpoints occurring in this argument have
norm at most \(R+1\). Write \(s=s_{\mu_P}\), \(s'=s_{\mu_Q}\),
and \(d=s'-s\).

Suppose two still equal limit parts have leading exponent \(E\), and
write a prefix-retaining step as
\(\rho=P_0+\omega^a\), with \(P_0\ne0\) retaining that leading
exponent. The exact next tail is
\begin{equation}\label{eq:local-phase}
 P_0+M(\omega^a,X),\qquad X=b_{i+1}-s(E)-N(P_0).
\end{equation}
The growth margin covers the retained prefix and a positive remainder;
any state below the prefix is inferior, and a tail at least
\(\omega^a\) is outside the interval. This proves the formula.

For \(a=1\) the new tail is finite, so a change of \(X\) does not
change the limit part. If \(a\) is a nonzero limit, the leading
exponent of the new tail is \(M(a,X-1)\), strictly increasing with
\(X\) by saturation. For \(a=a_0+1\), the exact recursion is
\begin{equation}\label{eq:successor-tail}
 M(\omega^a,X)=\omega^{a_0}q+M(\omega^{a_0},r),
 \quad q=\lfloor X/v\rfloor,\quad r=X\bmod v,
 \quad v=1+N(a_0).
\end{equation}
In particular, at \(a=2\) two different budgets have the same limit
part precisely when their difference is one and the larger budget is odd.

\begin{lemma}[A retained thick block forces separation]\label{lem:thick}
Suppose two undiverged relative descents process a common block
\(\omega^a q\), where \(a\ge2\) is fixed, \(N(a)\le R\), and
\(q>R+2\), while their whole-state leading exponent, and hence their
two distinct shifts \(c<c'\), remain fixed. Their limit parts diverge
before this block is processed.
\end{lemma}
\begin{proof}
If a lower remainder already creates a divergence, the conclusion holds.
Otherwise inspect the common next top visit
\(\rho=P_0+\omega^a(q-1)\). In each absolute descent the preceding
top visit lies at least \(\omega^2\) above this target. Both target
funnels are affordable. The long-entry residue formula, applied modulo
\(2M\) at the common preceding norm, gives
\[
 \left\lceil\frac{c+N(\rho)}2\right\rceil
 \equiv
 \left\lceil\frac{c'+N(\rho)}2\right\rceil\pmod M.
\]
The difference of these rounded quantities is smaller than \(M\),
so they are equal. It follows that \(c'=c+1\) and
\(c+N(\rho)\) is odd. At the next event the common budget \(u\)
satisfies, by the modular decrement,
\[
 u\equiv c+N(\rho)-1\pmod{2M}.
\]
Peel one unit, leaving \(P^-=P_0+\omega^a(q-2)\). The large
coefficient ensures that this prefix retains the original leading
exponent and hence the fixed shifts. The two local tail budgets are
\(X=u-c-N(P^-)\) and \(X-1\), with
\[
 X\equiv N(a)\pmod{2M}.
\]
If \(a\) is a limit, saturation makes the next leading exponents
different. If \(a=a_0+1\), then \(N(a)=1+N(a_0)=v\), and
\(v\mid M\). Thus \(v\mid X\). Formula
\eqref{eq:successor-tail} gives coefficients \(X/v\) and \(X/v-1\)
of the infinite \(\omega^{a_0}\)-block. The limit parts differ.
\end{proof}

This proof uses the moving residue of our control. A proof assuming that
every updated budget has one fixed residue would not apply to
\eqref{eq:factorial-control} or \eqref{eq:near-control}.

\begin{proposition}[An effective equal-gap certificate]\label{prop:equal-gap}
For two fixed equal-gap intervals one can compute a threshold and decide
whether their entry maps are identical on that tail or have one fixed
strict order there.
\end{proposition}
\begin{proof}
Here is the complete finite calculation. Let \(E=e(\delta)\).
During the initial phase, each nonleading term of \(\delta\) is
peeled one unit at a time, followed by all but one unit of its leading
term. The generated smaller blocks are processed before the next
original unit is peeled. Call a peeled unit sensitive if its exponent
is at least two.

If \(d(E)=0\), all prefix-retaining steps agree exactly. If
\(d(E)\ne0\), a sensitive limit exponent gives immediate divergence
by \eqref{eq:local-phase}. A successor exponent at least three either
gives immediate divergence or creates a common thick block, to which
Lemma~\ref{lem:thick} applies. Only exponent two can fail to separate.
Its coincidence condition is a shift difference of one and an odd larger
local budget. Consecutive original \(\omega^2\)-units have prefix
norms differing by three, whereas all updated budgets are even. Their
coincidence parities alternate, so two such units cannot both fail.
Thus the only nonempty sensitive initial phases which can fail to
separate have shapes
\[
 \delta=\omega^E+\omega^2+\omega c\quad(E\ge3),\qquad
 \delta=\omega^2\cdot2+\omega c,
\]
and their exact exception is
\[
 |d(E)|=1,\qquad \min\{s(E),s'(E)\}+N(P_0)\text{ odd},
\]
where \(P_0\) is the retained prefix at that one sensitive step.
If there is no sensitive step, the initial phase also fails to separate.
Every initial divergence has sign \(d(E)\) for \(E_P-E_Q\), by
candidate inclusion and Lemma~\ref{lem:propagation}.

If the initial phase has not separated, refine jointly the finitely many
breakpoint partitions of \(s\) and \(s'\). Search for the highest
piece on \(2\le e<E\) on which \(d(e)\ne0\). This is a finite
Cantor-normal-form calculation. An untruncated piece is right-closed,
so use its included right endpoint below \(E\). Only the top piece
clipped by the excluded endpoint \(E\) needs a different choice:
if \(E\) is a successor, use its predecessor when this lies in the piece;
if \(E\) is a limit and \(a\) is the lower edge, use the fixed exponent
\(a+1<E\). Discard an empty clipped piece and continue with the next one.

To justify that this finite piece determines the descent, note first
that phase exponents strictly decrease. A successor exponent drops to
its predecessor. A jump from a limit exponent cannot skip a fixed
\(\zeta\) of norm at most \(R+1\), since
\(\omega^{\zeta+1}\) is an affordable candidate. If the highest
nonzero piece has a maximum \(\beta\), the two descents agree until
that piece is reached. They enter it by a drop from \(\beta+1\):
a jump to the bounded \(\beta\) would have the larger affordable
candidate \(\omega^{\beta+1}\). The drop produces a common block
\(\omega^\beta q\) with \(q>R+2\), unless it has already diverged.
Lemma~\ref{lem:thick} then gives sign \(d(\beta)\).

A highest piece without a maximum is necessarily the top truncated
piece \((a,E)\), where \(E\) is a limit. On leaving the initial
pure \(\omega^E\)-state, a fixed exponent above \(a\) is affordable,
so both winning exponents lie within that piece. There the shifts are
constant, and the two exponents are
\(M(E,b_{i+1}-c-1)\) and \(M(E,b_{i+1}-c'-1)\).
Saturation separates them, with sign determined by the smaller shift.

If no such piece exists, all phases at exponents at least two finish
with identical limit parts and common budget index. For \(E\ge2\)
the final pure \(\omega^2\)-phase has next coefficients
\[
 b_{i+1}/2-\lceil s(1)/2\rceil,\qquad
 b_{i+1}/2-\lceil s'(1)/2\rceil.
\]
Their difference decides the order. Equal coefficients give simultaneous
termination. For \(E=1\), the original \(\omega\)-coefficient
already gives the same exact count in both descents.

The algorithm therefore uses, in order: the nonexceptional initial
phase; the highest nonzero profile piece on \([2,E)\); the difference
\(\lceil s'(1)/2\rceil-\lceil s(1)/2\rceil\); and otherwise
identity. All margins used above are covered by the common
tail. Equal stopping times give equal entry norms by
\eqref{eq:entry-orbit}; first limit-part divergence gives strict order
by propagation. This proves both completeness and effectiveness.
\end{proof}

\subsection{Cocycles, inverse thresholds, and collisions}

For fixed \(0\le\mu<\kappa<\lambda\), entry through \(\kappa\) requires
its affordability. On a tail satisfying all three endpoint norm bounds,
literal continuation gives
\begin{equation}\label{eq:cocycle}
 \begin{split}
 E_{\lambda,\mu}(t)&=
 E_{\kappa,\mu}(E_{\lambda,\kappa}(t)),\\
 \tau_{\lambda,\mu}(t)&=\tau_{\lambda,\kappa}(t)+
 \tau_{\kappa,\mu}(E_{\lambda,\kappa}(t)).
 \end{split}
\end{equation}
Genuine definedness of an exit alone is insufficient: at \(\lambda=\omega^2\),
\(t=3\), and \(B=f(3)\), the first predecessor is
\(\omega\cdot(B/2)\), which skips the fixed target
\(\kappa=\omega\cdot(B/2+1)\). A uniform sufficient guard for the
cocycle is \(t\ge\max\{2,N(\lambda),N(\kappa),N(\mu)\}\).

\begin{proposition}[Exact inverse-threshold table]\label{prop:threshold}
Given proper intervals \(P,Q\), comparisons of \(E_P(x)\) and
\(E_Q(y)\) have effective finite descriptions using comparisons of raw
inputs with fixed entry maps on the other input. Equality cases are
decided by the same description.
\end{proposition}
\begin{proof}
For equal gaps, Lemma~\ref{lem:gap-input} handles unequal inputs and
Proposition~\ref{prop:equal-gap} handles equal inputs. Suppose
\(\delta_Q<\delta_P\). Put
\[
 \kappa=\mu_P+\delta_Q,\quad
 \widetilde Q=(\kappa,\mu_P),\quad
 H=(\lambda_P,\kappa),\quad h=E_H(x).
\]
All these endpoints are effective fixed Cantor normal forms, and
\(\mu_P<\kappa<\lambda_P\). Impose their norm bounds and the relevant
comparison thresholds. The cocycle gives
\(E_P(x)=E_{\widetilde Q}(h)\).
Let \(\sigma\) be the eventual sign of
\(E_{\widetilde Q}-E_Q\) at a common input, supplied by the equal-gap
certificate. Input comparison brackets
\[
 E_Q(h-1)<E_{\widetilde Q}(h)<E_Q(h+1).
\]
Enlarge the tail so \(h-1\) is within all required domains. The exact
comparison of \(E_Q(y)\) with \(E_P(x)\) is consequently
\begin{center}\small
\begin{tabular}{@{}llll@{}}\toprule
\(\sigma\)& smaller & equal & larger\\\midrule
\(0\)&\(y<h\)&\(y=h\)&\(y>h\)\\
\(1\)&\(y\le h\)&never&\(y\ge h+1\)\\
\(-1\)&\(y\le h-1\)&never&\(y\ge h\)\\\bottomrule
\end{tabular}
\end{center}
Swap roles for the other gap order. If both original targets are zero,
\(\widetilde Q=Q\), so the middle certificate is identity.
Fixed input translations remain explicit in this table.

Every proper entry value belongs to the range of \(f\). Two distinct
such values \(U<V\) satisfy \(V\ge2^{U+1}\), by adjacent growth.
On a larger effective tail any fixed output offsets are therefore
irrelevant except at a collision, where their comparison is an ordinary
integer test. This also supplies the asserted descriptions with fixed
output offsets.
\end{proof}

The construction introduces the interval \(H\) and also the certificate
partner \(\widetilde Q\). Both must be retained when computing the
finite collection of source norms and thresholds. No inverse function
is added as a new primitive.

\section{Deciding closed blocks of collapse arithmetic}\label{sec:arithmetic}

We now combine order comparison with the modular tables. A finite circuit
may use natural inputs, natural-valued affine gates, fixed iterates of
\(f\), entry norms, interval counts, and fixed-head collapses. Each gate
and call carries its true domain; in particular an affine gate's output
is required to be natural. Signed linear combinations remain unrestricted
in the final inequalities, equalities and congruences. The total shifted function
\(n\mapsto J_\lambda(N(\lambda)+n)\) is a convenient way to avoid
partial raw inputs. Alternatively, an explicitly specified computable
finite extension below \(N(\lambda)\) can be handled as finite cases.

\begin{theorem}[Closed-block decision]\label{thm:arithmetic}
For either control of Section~\ref{sec:controls}, closed existential
and closed universal Boolean combinations of linear inequalities,
equalities and congruences over finite circuits of fixed-head collapses
are decidable. The procedure preserves shared inputs and shared circuit
nodes. The fixed heads may be arbitrary ordinals below \(\varepsilon_0\).
\end{theorem}

The descent notation below uses nonzero limit heads. This loses no fixed
head in the theorem: write any infinite \(\alpha\) uniquely as
\(\lambda+k\), with \(\lambda\) a nonzero limit and \(k<\omega\),
and absorb \(k\) into the finite input of
\(\psi_f(\alpha+u)=\psi_f(\lambda+(k+u))\).
For finite \(\alpha\), the value is simply \(\alpha+u\).

The proof is constructive, but makes no useful running-time claim. Its
two termination measures are the number of raw coordinates and a finite
charge on the active count occurrences.

\subsection{The finite matrix language}

Flatten a finite circuit once, assigning one natural coordinate to each
node and conjoining its exact graph equation. A collapse gate is
\[
 y=Z_\lambda(t)+D_\lambda(t),\qquad t\ge N(\lambda).
\]
Affine gates give affine equalities. Every valuation in the circuit's
original domain has exactly one extension to these graph coordinates.
This remains true when the
original output test is negated. All graph equations stay conjoined
outside that test.

After flattening, every nonlinear function has one raw coordinate as its
argument. Fixed argument shifts are represented by affine aliases. The
ordinary nonlinear value fields are
\[
 E_P(x),\quad f^a(E_P(x))\ (a\ge1),\quad f^a(x)\ (a\ge1),
\]
and the count fields are \(\tau_P(x)\). Coefficients in rows and
congruences are fixed integers. An uniterated raw coordinate is linear,
not a separated value field. The inner normalization introduces no new
quantified coordinates and no function of a newly constructed linear
combination.

\subsection{Comparison of compound values}

For \(P=(\lambda,\mu)\), put \(\nu=\mu+\omega\). On a tail
satisfying the bound \(x\ge N(\nu)\), define \(K_P(x)=x\) if \(\lambda=\nu\),
and \(K_P(x)=E_{\lambda,\nu}(x)\) otherwise. The last interval
has exactly one event, so the cocycle gives
\begin{equation}\label{eq:final-peel}
 E_P(x)=f(K_P(x)).
\end{equation}

\begin{lemma}[Comparison preserving input order]\label{lem:field-comparison}
On each fixed order of the raw coordinates, every comparison of ordinary
value fields has an effective finite Boolean form in which the greatest
raw coordinate occurs only linearly. Nonlinear parameters in that form
are fields on strictly smaller coordinates.
\end{lemma}
\begin{proof}
Cancel common outer applications of \(f\) by injectivity. If an explicit
iterate remains on one side and a primitive entry on the other, peel one
last \(f\) from the latter by \eqref{eq:final-peel} and cancel again.
Each such step decreases the remaining fixed iterate count. If a primitive
becomes the identity convention, retain the resulting comparison with a
raw coordinate. Thus finitely many steps leave two primitive entries, or
a field compared with a raw coordinate. Only finitely many endpoint
advances \(\mu+\omega j\) occur.

Identify equal raw coordinates in the chosen order branch. For primitive
entries at \(x<y\), Lemma~\ref{lem:gap-input} fixes the answer if the
larger-input map has at least the larger gap. Otherwise the inverse
threshold table makes it a comparison of \(y\) with an entry at \(x\),
with its exact equality condition. At equal inputs use the gap and
equal-gap certificates. A field at the larger input is greater than that
input and hence greater than a smaller bare coordinate. A larger bare
coordinate compared with a lower-input field is already linear in the
larger coordinate. Identity versus identity is raw order. These cases
are exhaustive, and all new source and certificate domains are retained.
\end{proof}

\subsection{Annotating a count by its full leading value}

A short interval \(\lambda=\mu+\omega k\) has constant count \(k\).
For every other proper interval, put
\(\eta=\mu+\omega^2\). Let \(P'=(\lambda,\eta)\), using the
identity convention if \(\lambda=\eta\). Its actual count satisfies
\begin{equation}\label{eq:count-core}
 \tau_P(x)=\tfrac12 C_P(x)+\tau_{P'}(x)+\kappa_\mu,
 \quad C_P(x)=f(E_{P'}(x)),\quad
 \kappa_\mu=1-\lceil N(\mu)/2\rceil.
\end{equation}
This is the exact \(\omega^2\)-funnel formula, including the
\(\omega\)-tail of \(\mu\) absorbed by \(\eta\).
For \(\mu=0\) it reads \(D_{\omega^2}(q)=f(q)/2+1\), after the
preceding entry prefix. In every case, if \(q=E_{P'}(x)\) and
\(r=\tau_{P'}(x)\), then on the genuine tail
\begin{equation}\label{eq:core-error}
 q\ge x,\qquad C_P(x)\ge f(x),\qquad
 0\le r\le q\le\log_2 C_P(x).
\end{equation}

We attach \((C_P,P',\kappa_\mu)\) as an annotation to the whole count
occurrence. We do not expand it into a new independent half-value and
an active prefix count. In a class sharing the full value \(C\), ordinary
coefficient \(b\) and count coefficients \(a_i\) have combined leading
coefficient \(b+\frac12\sum_i a_i\). Its zero test is the integer test
\(2b+\sum_i a_i=0\). Treating \(C\) and \(C/2\) as separated
scales would give an incorrect test.

\begin{lemma}[Exact cancellation and transfer]\label{lem:count-transfer}
A class with zero combined leading coefficient can be replaced exactly by
its prefix counts and constants. Equal-input prefixes merge. Prefixes on
distinct inputs can be transferred to smaller input ranks whenever their
cores coincide.
\end{lemma}
\begin{proof}
For a zero class, \eqref{eq:count-core} gives the exact replacement
\begin{equation}\label{eq:zero-replacement}
 bC+\sum_i a_i\tau_{P_i}(x_i)
   =\sum_i a_i\tau_{P_i'}(x_i)+\sum_i a_i\kappa_i.
\end{equation}
All ordinary value occurrences in the class disappear. Identity and
short prefixes become constants. If two remaining prefixes have the
same input \(x\), equality of their cores gives equality of their entry
norms. These are \(f^r(x)\) and \(f^s(x)\); strict increase above the
identity gives \(r=s\). Their counts therefore merge on this guard.

Now take two remaining prefixes at \(x<y\), with
\(E_{P'}(y)=E_{Q'}(x)\). They are proper after the preceding removals.
Input comparison forbids \(\delta_{P'}\ge\delta_{Q'}\).
The complete collision table therefore forces its equality certificate
and
\[
 y=E_H(x),\qquad
 H=(\lambda_{Q'},\mu_{Q'}+\delta_{P'}).
\]
The actual \(Q'\)-cocycle factors through \(H\). At \(y\) its
remaining same-gap entry has the same terminal norm as \(P'\), hence
the same count by injectivity of iteration. Consequently
\begin{equation}\label{eq:rank-transfer}
 \tau_{P'}(y)=\tau_{Q'}(x)-\tau_H(x).
\end{equation}
This is an integer identity on the actual collision guard, even with
signed coefficients. A nonequality certificate makes the collision
branch impossible. Both terms on the right have smaller raw input rank.
Choose the smallest surviving input and transfer every higher one to it.
All intermediate target norm bounds and comparison domains are imposed.
\end{proof}

\subsection{A finite charge and the order normalization loop}

Fix a strict order of the \(r\) distinct raw coordinates, giving ranks
\(0,\ldots,r-1\). For one row put
\begin{equation}\label{eq:finite-charge}
 \mathcal M=\sum_{\text{active count occurrences}}3^{\operatorname{rank}(x)},
 \qquad
 \mathcal V=\#\{\text{ordinary nonlinear value occurrences}\}.
\end{equation}
Combine identical syntactic fields and remove zero coefficients.
Annotations add no ordinary occurrences. Pattern guards are kept
separately; Lemma~\ref{lem:field-comparison} already makes them linear
in the greatest coordinate.

\begin{lemma}[Finite inner normalization]\label{lem:inner-termination}
Every exact zero-class replacement strictly decreases
\(\mathcal M+\mathcal V\), and neither summand increases.
\end{lemma}
\begin{proof}
An ended prefix removes a count occurrence. Merging two prefixes at
one input replaces at least two equal-rank occurrences by at most one.
A transferred occurrence of rank \(j\) creates at most two lower-rank
occurrences, of total charge at most
\(2\cdot3^{j-1}<3^j\). If there is just one surviving prefix and
its count charge has not decreased, the equation \(2b+a=0\) with
\(a\ne0\) forces \(b\ne0\); an ordinary value occurrence has
been consumed. A class consisting only of ordinary values also decreases
\(\mathcal V\). Every zero class with two or more count occurrences
has a repeated input, distinct inputs, or an ended prefix, so one of
the preceding cases applies. Formula \eqref{eq:zero-replacement} creates
no ordinary values. This proves the assertion.
\end{proof}

Let \(z\) be the greatest raw coordinate. Normalize an order atom by
the following finite procedure. Remove short counts and zero coefficients.
If \(z\) has no remaining nonlinear occurrence, retain the entire row
as a linear row in \(z\), with its correlated lower-input fields as
parameters, and stop. Otherwise annotate all remaining counts, branch on
the weak order of their full cores and the ordinary nonlinear fields,
and retain every exact pattern guard. Process the largest current class.
A zero combined coefficient is canceled by Lemma~\ref{lem:count-transfer},
after which the procedure repeats. At a nonzero coefficient it stops by
the following domination argument.

Because the initial exit check failed, a nonlinear field at \(z\)
remains, so the largest current full value \(C\) satisfies \(C\ge f(z)\).
Every raw coordinate is at most \(z\le\log_2 C\). Distinct ordinary
full values are separated by adjacent growth; every smaller one is at
most \(\log_2 C\). Count errors satisfy \eqref{eq:core-error}.
The entire row is consequently
\[
 \beta C+E,\qquad 0\ne\beta\in\tfrac12\mathbb Z,
 \qquad |E|\le K\log_2 C+K_0,
\]
for effective fixed integers \(K,K_0\). An elementary cutoff makes
\(|E|<C/2\), so its sign is the sign of \(\beta\), and equality is
false. Previously canceled errors have not been dropped: they were
replaced exactly by active prefixes before this new largest class was
chosen. If those cancellations removed all nonlinear occurrences of
\(z\), the exit check instead retained its exact linear row.

The charge lemma bounds the depth of every inner branch by the initial
\(\mathcal M+\mathcal V\). Each node has finite branching. Thus the
entire symbolic normalization tree can be generated before numerical
tails are chosen. It contains finitely many interval, certificate,
peeling and transfer labels, and finitely many coefficient bounds.

\subsection{Residues, finite strips, and one-variable projection}

Proposition~\ref{prop:entry-residues} converts every current field
congruence modulo \(m\) into a finite Presburger condition on its
original raw input modulo \(2m\). For example a long count has residue
\(p(x)-c_\mu\), while a short entry has residue \(2p(x)-2k\).
These are tables in \(x\bmod2m\). All occurrences of the same \(x\)
use one shared residue. An optional explicit half-field can either use
the corresponding table or its exact natural graph \(2h=f(q)\).
The count annotations themselves introduce no such half-field atom.

For a matrix in \(n\) coordinates, first generate all finite inner
normalization trees and all current modular tables. Choose a common
cutoff \(B\) above their source norms, affordability and comparison
thresholds, prefix cutoffs, and domination margins. This maximum is
effective because the symbolic closure is finite. It does not require
finding an unknown first divergence or computing a descent length.

If some coordinate is below \(B\), branch over each of its finitely
many permitted values and substitute that same value throughout the
matrix. All fields on it are now closed exact calculations, with their
true definedness flags. The other coordinates stay unbounded: these
branches cover whole infinite strips, not merely a finite cube. Each
has one fewer coordinate and retains every earlier domain guard.

On the all-large branch, identify equal coordinates or choose a strict
base order. The order normalization and every pattern guard are linear
in the greatest coordinate \(z\). Current congruences are ordinary
Presburger congruences in the raw coordinates. Apply effective Presburger
elimination to \(\exists z\), treating all lower-input fields as
parameters \cite{cooper}. Its equivalence holds for every parameter
assignment; substituting the actual correlated values preserves it.
The resulting formula is in the same matrix language on \(n-1\)
coordinates: coefficients and moduli may change, but arguments of fields
do not change.

New moduli are handled at the next, smaller-dimensional stage, with its
own finite tails and complete small-input strips. There is no loop
raising cutoffs at the same dimension. At zero coordinates use the
exact closed evaluator. Induction on \(n\), together with the inner
charge, proves termination and exactness. A universal block is the
absence of an existential counterexample on the same original domain.
Unique circuit extensions justify the initial flattening. This proves
Theorem~\ref{thm:arithmetic}.

\subsection{Two checks on the arithmetic mechanism}

The identities
\[
 D_{\omega^2}(t)=f(t)/2+1,\qquad
 2D_{\omega^2}(t)-f(t)-1=1
\]
illustrate why full values, rather than literal half-values, must be
grouped. The leading coefficient cancels, leaving the exact constant.
For a transfer between distinct input ranks, put
\(\mu=\omega^2\), \(\lambda=\omega^2\cdot2\), and \(y=f(x)\).
The first move from \(\lambda+\omega\) reaches the \(\lambda\)-cut,
so
\[
 D_{\lambda+\omega}(x)=1+D_\lambda(y).
\]
The two long-count prefixes have equal cores. Their collision transfers
the higher-input count to the lower input with an exact difference of
one. These are symbolic identities for the genuine control; no smaller
numerical surrogate is used.

The theorem concerns closed blocks. It does not establish quantifier
elimination with retained free parameters in the same forward language.
Indeed the unary quantifier-free formulas of the retained matrix language
are eventually periodic, by the same order normalization and residue tables. The projection
\(\exists z\,[x=f(z)]\) is an infinite set with unbounded gaps and
is not eventually periodic. Thus a general projection theorem would
need a richer output language. This distinction does not obstruct the
closed counterexample sentences used for canonical comparison.

\section{The canonical comparison theorem}\label{sec:consequences}

Fix one computable admissible literal parsing policy \(\rho\). A
canonical assignment below means the zero-history assignment of the
literal input term under that policy, not the assignment of a normalized
term and not an arbitrary history-dependent assignment. Write
\(\can A_{f,\rho}\) for this vector.

\begin{theorem}[All ground free variables]\label{thm:ground}
For either \(f_*\) or a fixed \(f_b\), universal weak comparison,
Wilken--Weiermann strict comparison, and universal equality of canonical
assignments are decidable for original G\"odel T terms whose free variables
all have ground type. There is no bound on internal types, bound-variable
types, result types, recursor instances, or finite nesting.
\end{theorem}
\begin{proof}
The generated-source theorem gives, for the two supplied terms,
\[
 \can A_i=a_i,\quad\can B_i=b_i\quad(i>0),\qquad
 \can A_0=E_A(\mathbf x),\quad\can B_0=E_B(\mathbf x),
\]
where the positive coordinates are closed computable ordinals and the
scalar expressions are finite circuits built from positive natural inputs,
constants, sums, fixed natural multiples, and
\(u\mapsto\psi_f(\lambda+u)\) with fixed nonzero limit heads
\(\lambda<\varepsilon_0\). Repeated names have one shared coordinate.
Every positive integer is an admissible ground value, represented by a
finite sum of ones, so these inputs range independently over
\(\N_{>0}\).

First compute the finite positive-coordinate comparisons, using zero
padding beyond the assigned level. If any required one fails, the
vector comparison fails. Otherwise universal weak comparison is exactly
the absence of a solution to
\[
 \exists\mathbf x\in\N_{>0}^{\,r}\quad E_A(\mathbf x)>E_B(\mathbf x).
\]
For strict comparison the counterexample relation is \(\ge\), and
for equality it is \(\ne\). Flatten the scalar circuits using their
unique graph coordinates and true norm guards. Every graph coordinate
and every original input belongs to the same existential block.
Theorem~\ref{thm:arithmetic} decides this closed sentence. Internal
abstractions have already been handled by the finite assignment recursion;
they do not create external quantifier alternation. This proves all
three claims.
\end{proof}

The result is a forward translation theorem. It does not require a typed
source realizing every auxiliary interval created by arithmetic
elimination. Those intervals are legitimate arithmetic objects with the
explicit definitions proved above.

\begin{lemma}[Fixed-head tail growth]\label{lem:tail-growth}
For a nonzero limit \(\lambda\), let
\(h_\lambda(u)=\psi_f(\lambda+u)\). Then
\[
 h_\lambda(u+1)>2^{h_\lambda(u)}.
\]
\end{lemma}
\begin{proof}
Increasing the initial norm by one cannot shorten its nonfinite descent:
monotonicity of maximal predecessors, followed by taking limit parts,
propagates the weak comparison. Write the smaller terminal state as
\(Z=f^D(t)\). The larger terminal state is at least \(f^D(t+1)\).
At the last iterate, the preceding integer arguments differ by at least
one, so \eqref{eq:growth} gives \(f^D(t+1)\ge2^{2^Z}\).
Since \(D<Z\), the smaller collapse is \(Z+D<2Z\le2^Z\).
The terminal decomposition proves the claimed inequality.
\end{proof}

As an actual-source instance, with the native recursor argument order put
\(H=R^o0(D^o0):o\to o\). Its first positive coordinate is
\(\omega\cdot2\). Thus the application \(HX\) invokes the head
\(\omega^2\cdot2\) at the genuine raw norm
\(N(\omega^2\cdot2)+c_H+x_0\), where \(c_H\) is its actual closed
scalar. The complete-count table applies along this input ray without
replacing the source by an arbitrary ordinal-valued primitive.

\section{Undecidability with functional inputs}\label{sec:functional-hardness}

The same controls that give the ground-input decision theorem support
an undecidable comparison problem once free variables of type $o\to o$
are allowed. The reduction uses the growth estimates, not any modular
property of the control. The final corollary records the same negative
result for the unchanged standard control.

The published undecidable problem used here is solvability of a finite
system of linear equations
\[
 \sum_i a_{ji}(X)p_i(X)-g_j(X)=0,
 \qquad a_{ji},g_j\in\mathbb Z[X],\quad p_i\in\mathbb N[X].
\]
It is linear in the unknown polynomials, not a system of arbitrary
Diophantine products of unknowns.  Narendran's theorem gives this precise polynomial-semiring boundary
\cite{narendran}.
Zero polynomials are allowed.  We use undecidability, without importing a
stronger many-one completeness classification of that source problem.

For ground sources define the eraser $e(U,V)=(\lambda u:o.U)V$,
with a fresh unused binder. Its literal scalar is
\[
 e(U,V)_0=g(3U_0+V_0+2),\qquad g(n)=\psi_f(\omega+n)=f(n+2)+1.
\]
The function $g$ is strictly increasing with gaps at least two.

Put $\tau=o\to o$.  For actual sources $A,B:\tau$ define
\begin{align*}
 \mathsf P(A,B)&=(\lambda z:\tau.A)B,\qquad
 \mathsf H(A)=R^o0(\lambda n:o.A),\\
 \mathsf S(A,B)&=(\lambda z:(o\to\tau).A)(\lambda n:o.B).
\end{align*}
The original argument order is
$R^\rho:o\to(o\to\rho\to\rho)\to\rho\to\rho$:
its first argument $0:o$ is the recursion test, not its base value.
Thus $R^\rho0(\lambda n:o.A)$ has type $\rho\to\rho$
when $A:\rho\to\rho$, also at the higher type used below.
The displayed bound names are unused and fresh. These literal identities
are uniform over admissible parsing policies: their unused-binder partials
do not inspect a rebracketing inside an active variable occurrence.  Writing
$a_i=\can A_i$, $b_i=\can B_i$, the literal canonical clauses give
\begin{align*}
 \can{\mathsf P(A,B)}_1&=8a_1+4b_1+8,\\
 \can{\mathsf S(A,B)}_1&=8a_1+8b_1+16,\\
 \can{\mathsf S(A,B)}_0
   &=\psi_f\bigl(\omega\otimes(8a_1+8b_1+16)+9a_0+5b_0+13\bigr),\\
 \can{\mathsf H(A)}_1&=2\omega\otimes a_1+3\omega.
\end{align*}
For example the unused $\tau$ binder has vector
$(9a_0+8,2a_1+2,2)$, and its box with $B$ gives the first formula.
For the second binder its scalar is again $9a_0+8$ and its first field
$2a_1+2$; the ground lambda on $B$ has vector
$(5b_0+4,2b_1+2)$, giving both $\mathsf S$ formulas.
The recursor prime followed by its two displayed arguments gives
the $\mathsf H$ formula.  These calculations use actual source syntax;
they do not introduce an arbitrary ordinal-valued primitive.

\begin{lemma}[Signed polynomial pairs]\label{lem:pair-compiler}
Every signed linear expression $D$ in ordinal unknowns with coefficients
in $\mathbb Z[\omega]$ has effectively constructed actual $A_D,B_D:\tau$
whose first-coordinate difference is $2^kD$ for a finite $k$.
All free functional variables have type $\tau$, and internal types have
level at most two.  Differences are taken only in the signed Cantor-monomial
group; each of the two source coordinates remains a positive ordinal.
\end{lemma}
\begin{proof}
Let $C=\lambda u:o.u$, with $C_1=1$.  The first fields of
$\mathsf P(C,C)$ and $\mathsf S(C,C)$ are 20 and 32.  Take sixteen pairs:
eight copies of $(\mathsf S(C,C),\mathsf P(C,C))$ of difference 12,
five of $(C,\mathsf P(C,C))$ of difference $-19$, and three of $(C,C)$.
Combine them componentwise in a balanced $\mathsf S$ tree of height four.
Constants cancel in every paired combination, so the resulting difference is
\[
 8^4(8\cdot12-5\cdot19)=2^{12}.
\]
For a free $X:\tau$, the pair
$(\mathsf P(X,C),\mathsf P(C,C))$ has difference $8(x_1-1)$.
Three multiplications by eight followed by a paired $\mathsf S$ combination
with the constant pair give $2^{15}x_1$.

The pair operation $(A,B)\mapsto(\mathsf S(A,C),\mathsf S(B,C))$
multiplies a difference by eight, and
$(A,B)\mapsto(\mathsf S(A,A),\mathsf S(B,B))$ multiplies it by sixteen.
Applying $\mathsf H$ to both entries multiplies a difference by $2\omega$;
swapping entries negates it.  Repeated addition constructs every finite
integer coefficient.  To add pairs with scale exponents $k_1,k_2$, choose
$K\ge\max(k_1,k_2)+6$.  Every integer at least six is $3u+4v$ with
$u,v\ge0$, so the two scale operations align both differences at $2^K$.
A paired $\mathsf S$ then gives their sum with scale $2^{K+3}$.
Induction on the finitely many signed monomials proves the assertion.
The displayed source combinators introduce no internal level above two.
\end{proof}

The two scalar coordinates of an equation pair need not agree when their
positive coordinates agree.  We therefore construct that alignment explicitly.

\begin{lemma}[Alignment and a full-vector detector]\label{lem:functional-detector}
At each fixed admissible valuation of the original inputs, the condition
$\can A_1=\can B_1$ is equivalent to the existence of two fresh
positive ground inputs making the full canonical vectors of
\[
 A'=(\lambda t:o.A)z,\qquad B'=(\lambda t:o.B)w
\]
equal.  Full-vector equality of any actual pair of type $\tau$ can in turn
be detected by ground sources $L,R$ satisfying $L_0\ge R_0$, with equality
exactly at vector equality and otherwise gap at least two.
\end{lemma}
\begin{proof}
The aligned first fields are $2a_1+2$ and $2b_1+2$, and their scalars are
\[
 \psi_f(\omega\otimes(2a_1+2)+5a_0+4+z_0),\qquad
 \psi_f(\omega\otimes(2b_1+2)+5b_0+4+w_0).
\]
Equal vectors imply $a_1=b_1$.  If this equality holds, choose
$z_0=5b_0+1$, $w_0=5a_0+1$.  Both tails and both vectors then agree.
These choices are made after the original valuation, so are not circular.

For the detector, now write $A,B$ for the pair being tested, and form
\[
 X=\mathsf S(A,A),\quad Y=\mathsf S(B,B),\quad
 Z=\mathsf S(A,B),\quad W=\mathsf S(B,A).
\]
Make the actual shared ground payload
$Q=W(Z(Y(X0)))$, and put $q=Q_0$.
The unconditional finite-tail lower bound for $\psi_f$ gives
\[
 q\ge1+x_0+y_0+z_0+w_0\ge5.
\]
Set $g_X=XQ$, and similarly $g_Y,g_Z,g_W$, and finally set
$L=e(g_X,g_Y)$, $R=e(g_Z,g_W)$.

Suppose first $a_1>b_1$.  Then $x_1>z_1=w_1$.  The ordinals
\[
 \alpha=\omega\otimes z_1+z_0+q+2,
 \qquad \beta=\omega\otimes x_1+x_0+q
\]
satisfy $\alpha<\beta$.  We must check a norm premise before comparing
their collapses.  Bounded norms of actual assignments and the published
natural-product norm inequality \cite[Propositions 2.2--2.3]{ww} give
\[
 \norm(\alpha)\le2\norm(z_1)+z_0+q+2
 \le3z_0+q+2\le4q+2
 \le2^{q+1}\le f(\norm(\beta)).
\]
Controlled strictness therefore gives $\psi_f(\beta)>\psi_f(\alpha)$.
Two consecutive fixed-limit tail-growth steps from Lemma~\ref{lem:tail-growth}
yield $g_{X,0}>4g_{Z,0}$.  The same calculation with $W$ gives
$g_{X,0}>4g_{W,0}$.  If $b_1>a_1$, use $Y$ in place of $X$.

If $a_1=b_1$ but $a_0\ne b_0$, the four $\mathsf S$ heads coincide.
Their scalar tails are
\[
 14a_0+13,\quad14b_0+13,\quad
 9a_0+5b_0+13,\quad9b_0+5a_0+13.
\]
The larger diagonal tail exceeds both cross tails by at least five.
Strict increase gives at least a two-unit scalar gap, so two fixed-tail
growth steps in the subsequent applications give again a diagonal $g$
value greater than four times each cross value.

In either unequal case,
$3g_{X,0}+g_{Y,0}>3g_{Z,0}+g_{W,0}$: even when the larger diagonal
is the one with coefficient one, its fourfold domination suffices.
The final eraser is increasing and has gaps at least two.
Equal input vectors give four equal $\mathsf S$ vectors, hence equality.
This proves the claimed detector on every admissible valuation, not only
on the specially chosen alignment witnesses.
\end{proof}

\begin{theorem}[Functional-input undecidability]
\label{thm:functional-hardness}
Full universal canonical weak comparison and strict comparison are
undecidable already for ground output, free-variable levels at most one,
and internal levels at most two.  The reduction at this level uses finitely
many comparison queries.  A computable single-query reduction from the
same polynomial-system unsatisfiability problem is available with internal
level at most three.
\end{theorem}
\begin{proof}
Let $\pi$ retain exactly the Cantor monomials with finite exponents.
It maps ordinals below $\varepsilon_0$ to $\mathbb N[\omega]$.
It preserves natural sum and multiplication: a natural sum of exponents
is finite exactly when both exponents are finite.  It consequently extends
to the signed Cantor-monomial group and respects multiplication by fixed
$\mathbb Z[\omega]$ coefficients.

For the level-two assertion, enumerate the finitely many possible sets of
zero polynomial unknowns, substitute zero for those, and require the others
to be positive ordinal first coordinates.  Compile the resulting equations
by Lemma~\ref{lem:pair-compiler}, align them separately, and apply
Lemma~\ref{lem:functional-detector}.  Aggregate the nonnegative ground
detectors in identically shaped eraser trees.  Root equality is exactly
simultaneous equation satisfaction.  Its absence is one universal strict
comparison; the conversion
\[
 \can{KR}_0\le\can{IL}_0\quad\Longleftrightarrow\quad
 g(5+R_0)\le g(3+L_0)\quad\Longleftrightarrow\quad L_0-R_0\ge2
\]
 gives the same weak query. Here $K=\lambda z:o.0$ and
$I=\lambda z:o.z$, with canonical vectors $(5,1)$ and $(3,1)$.
An empty polynomial system is first augmented by the tautological equation
$0=0$, so the same unsatisfiability-query convention applies.

A polynomial solution lifts in the matching zero pattern: each remaining
nonzero polynomial is an admissible positive first field, after choosing its
zeroth field sufficiently large for its norm.  Conversely any positive-ordinal
solution projects by $\pi$ to a polynomial solution of the original system.
Some projected unknowns may become zero, which is permitted in that system.
Thus an algorithm for either comparison would decide polynomial solvability
using the finitely many patterns, contradicting the published theorem.

For a single query choose instead a fixed positive ordinal $\Omega$ with
$\pi(\Omega)=0$, and transform each original equation to
\[
 \sum_i a_i(\omega)x_i-g(\omega)
                    -\left(\sum_i a_i(\omega)\right)\Omega=0.
\]
A polynomial solution $u_i$ lifts by $x_i=u_i\oplus\Omega>0$.
Conversely projection of any positive-ordinal solution solves the original
equations.  The transformation preserves simultaneous use of each unknown.

We must realize the extra constant in actual source syntax.  Put
$\sigma=\tau\to\tau$, $C_2=\lambda u:\tau.u$, and
$\mathsf H_2(A)=R^\tau0(\lambda n:o.A)$.  The canonical clauses give
$(C_2)_2=2$ and $(\mathsf H_2(A))_2=\omega+2\omega\otimes A_2$.
The closed $H=\mathsf H_2(\mathsf H_2(C_2))$ has
\[
 H_2=10\omega^2+\omega,
 \qquad\Omega=2^{10\omega^2+\omega}=\omega^{10\omega+1},
 \qquad\pi(\Omega)=0.
\]
Apply this same actual $H$ to the two entries of the constant pair of
Lemma~\ref{lem:pair-compiler}.  The box gives
$\can{HA}_1=\Omega\otimes(H_1+A_1)$, so their difference is
$2^{12}\Omega$.  Add this as a second closed seed to the pair compiler.
Only fixed $\omega$ shifts and signed sums of that seed are needed.
Compile the one transformed system and apply the same alignment and detector.
The closed $R^\tau$ raises the internal level bound to three, while every
free variable still has level at most one.  This is the promised single-query
reduction.
\end{proof}

\begin{corollary}[Functional inputs at the unchanged standard control]
\label{cor:standard-functional-hardness}
The weak and strict undecidability statements of
Theorem~\ref{thm:functional-hardness} also hold for the unchanged control
\(G(n)=F_5(n+100)\), with the same finite-query internal level bound two
and single-query internal level bound three.
\end{corollary}
\begin{proof}
The canonical-source interface requires a strictly increasing control
majorizing \(G\), and the numerical separation used in the detector
requires the adjacent growth estimate \eqref{eq:growth}. The standard
control meets these premises. Indeed, with \(A=G(n)\), the hierarchy
calculation in Lemma~\ref{lem:control} gives
\[
 G(n+1)\ge F_4(A+1)\ge F_1(A+1)
             \ge 2^{2^{A+1}}\ge2^{2^{G(n)}}.
\]
It is primitive recursive and has the same eventual finite-level bound.
No count-residue or modular-potential statement occurs in the source
compiler, alignment, detector or undecidability reduction. Their literal
constructions and proofs therefore apply with \(f=G\) throughout.
\end{proof}



\section{Consequences and remaining questions}

The positive and negative halves use the same chosen control. Thus the
ground/function distinction is a decision-theoretic boundary in the
canonical source language, rather than a comparison between unrelated
assignments. The positive theorem also covers equality; no equality
undecidability assertion for functional inputs is needed or proved.

The prefix-preserving control shows that any prescribed finite set of
initial values of the published example is compatible with the explicit
residue construction and the positive decision theorem.
It supplies an existence result within the original permitted control
framework. It does not evaluate a previously unknown residue for the
unchanged example, and it does not identify the two resulting assignments.

The proof does not need the separate special-case algorithms for one
input, polynomial heads, balanced layers, or low internal type levels.
Those may remain useful for sharper bounds or for the unchanged control.
The present proof uses the all-head entry calculus and one finite count
normalization. An implementation, useful complexity bounds, a classification
at the unchanged control, and a full alternating first-order theory are
additional questions.

\section*{Preparation and verification status}

This manuscript was prepared through AI-assisted research directed by the
repository owner. No individual is assigned mathematical authorship by
this preparation record. The arguments printed here are hand proofs.
Earlier formal work checked several scalar, domain, and source-translation
lemmas and exposed the strict-inequality defects discussed in the text;
it does not constitute a proof-assistant verification of the new control,
entry comparison, or complete decision theorem. Build checks and separate
AI review are not independent human verification or a historical priority
assessment. All mathematical dependencies needed beyond the cited
published background are stated and proved in this manuscript.

\begin{thebibliography}{9}
\bibitem{bcw}
Wilfried Buchholz, Adam Cichon, and Andreas Weiermann.
A uniform approach to fundamental sequences and hierarchies.
\emph{Mathematical Logic Quarterly} \textbf{40}(2) (1994), 273--286.
\href{https://doi.org/10.1002/malq.19940400212}{doi:10.1002/malq.19940400212}.

\bibitem{cooper}
D. C. Cooper.
Theorem proving in arithmetic without multiplication.
In B. Meltzer and D. Michie, editors, \emph{Machine Intelligence 7},
pp.~91--99. Edinburgh University Press, 1972.
\url{https://www21.in.tum.de/teaching/logik/SS16/Exercises/Cooper.pdf}.

\bibitem{narendran}
Paliath Narendran.
Solving linear equations over polynomial semirings.
In \emph{Proceedings of the Eleventh Annual IEEE Symposium on Logic in
Computer Science}, pp.~466--472. IEEE Computer Society Press, 1996.
\url{https://lics.siglog.org/archive/1996/Narendran-Solvinglinearequati.html}.

\bibitem{ww}
Gunnar Wilken and Andreas Weiermann.
Derivation lengths classification of G\"odel's T extending Howard's assignment.
\emph{Logical Methods in Computer Science} \textbf{8}(1:19) (2012), 1--44.
\href{https://doi.org/10.2168/LMCS-8(1:19)2012}{doi:10.2168/LMCS-8(1:19)2012}.
\url{https://lmcs.episciences.org/1073}.
\end{thebibliography}

\end{document}
